\documentclass[11pt]{article}

\usepackage[T1]{fontenc}
\usepackage{lmodern}
\usepackage{microtype}
\usepackage{amsmath,amssymb,amsthm,mathtools}
\numberwithin{equation}{section}
\usepackage{enumitem}
\usepackage{array}
\usepackage{tabularx}
\usepackage{geometry}
\geometry{margin=1in}
\usepackage{cite}
\usepackage{hyperref}
\usepackage{url}
\usepackage{tikz}
\hypersetup{colorlinks=true,linkcolor=blue,citecolor=blue,urlcolor=blue}
\setlength{\emergencystretch}{3em}

\theoremstyle{plain}
\newtheorem{theorem}{Theorem}[section]
\newtheorem{lemma}[theorem]{Lemma}
\newtheorem{proposition}[theorem]{Proposition}
\newtheorem{corollary}[theorem]{Corollary}
\theoremstyle{definition}
\newtheorem{definition}[theorem]{Definition}
\theoremstyle{remark}
\newtheorem{remark}[theorem]{Remark}
\newtheorem{example}[theorem]{Example}
\newtheorem{openproblem}[theorem]{Open Problem}

\newcommand{\Wit}{\mathsf{W}}
\newcommand{\Cmcf}[2]{\mathcal{C}^{\mathrm{mcf}}_{#1,#2}}
\newcommand{\Cboundedtyping}[2]{\mathcal{C}^{\mathrm{mcf},\le #2}_{#1}}
\newcommand{\Calltyping}[1]{\mathcal{C}^{\mathrm{mcf},\mathrm{fin}}_{#1}}
\newcommand{\Ffrag}[2]{\mathcal{F}_{#1,#2}}
\newcommand{\D}{\mathcal{D}}
\newcommand{\out}{\mathrm{out}}
\newcommand{\tuple}[1]{\vec{#1}}
\newcommand{\blank}{\square}
\newcommand{\eps}{\varepsilon}
\newcommand{\pausesym}{\bot}
\newcommand{\fan}{\operatorname{fan}}
\newcommand{\YSub}[1]{\mathsf{YSub}(#1)}
\newcommand{\YSL}[2]{\mathsf{YSL}(#1,#2)}
\newcommand{\evalmap}{\mathsf{ev}}

\newcommand{\cutpos}[1]{%
  \raisebox{-0.7ex}{%
    \tikz[baseline=(char.base)]{%
      \node[draw,circle,inner sep=0.4pt,font=\scriptsize]
        (char) {\ensuremath{#1}};%
    }%
  }%
}

\title{Positive-Data Learning of Multiple Context-Free Languages under Fixed Finite-Monoid Typing}
\author{Takayuki Kuriyama\\
Independent Researcher, Tokyo, Japan\\
\texttt{growup.kuriyama@gmail.com}}
\date{}

\begin{document}
\maketitle

\begin{abstract}
Fix a fan-out bound $f$ and a homomorphism $h:\Sigma^*\to M$ into a finite
monoid.  We study positive-data learning of multiple context-free languages
under the bias that tuple components may be substituted only when their
corresponding $h$-values agree.  Starting from Yoshinaka's positive-data
reconstruction for multidimensional substitutability, we impose this finite
type guard on substitution links and prove exact reconstruction after a finite
target-dependent positive sample.  An uncapped conservative learner therefore
identifies every target satisfying the fixed fan-out and typing condition in
the limit; with a fixed rule-rank cap, updates are polynomial in the
accumulated sample.

The main quantitative distinction is by rule rank.  For every fixed finite
typing, rule-rank-one targets admit characteristic data of total length
polynomial in grammar size alone.  At arbitrary fixed rank, for reduced good
presentations, typings determined by fixed-length prefixes and suffixes admit
characteristic data polynomial in grammar size and rule thickness.  We combine a known non-branching obstruction with a boundary-typed wrapper
to obtain a $(1,1)$-boundary-typed non-context-free target of minimum fan-out
two, outside Yoshinaka's trivial-typing class, whose minimum rule rank is two
even when arbitrary finite fan-out is allowed.  Thus the higher-rank theorem
applies to a genuinely non-context-free target that lies beyond the rank-one
result at every finite fan-out.  Full yield-type refinement and empty-component elimination
can each increase rule thickness exponentially.

We also show that finite typing is not merely regular post-filtering and that
fan-out is part of the learning bias.  At every level $m\ge2$ there is a
rank-one target of minimum fan-out $m$ in a class typed by fixed first/last
symbol information but outside Yoshinaka's class, and at fan-out two another
finitely typed target cannot be represented as a Yoshinaka-substitutable
language intersected with any target-dependent regular language.
\end{abstract}

\medskip
\noindent\textbf{Keywords:} grammatical inference; multiple context-free grammars; positive data; finite-monoid typing; multidimensional substitutability; Gold identification; rule rank.
\medskip

\section{Introduction}
\label{sec:introduction}

Positive-data grammatical inference asks when a language can be recovered from
an enumeration of its positive examples alone \cite{Gold1967,Angluin1980}.
Distributional approaches make this possible for restricted grammar classes by
turning observed interchangeability of substrings or tuples into grammar
nonterminals.  For context-free grammars (CFGs) this includes substitutable languages
\cite{ClarkEyraud2007} and Yoshinaka's $(k,\ell)$ hierarchy
\cite{Yoshinaka2008}.  Yoshinaka extended the same idea to multiple
context-free grammars (MCFGs), where nonterminals generate tuples and
substitution is tested by multi-hole sentence contexts \cite{Yoshinaka2009,Yoshinaka2011}.
MCFGs form a standard mildly context-sensitive formalism and can directly
represent discontinuous and cross-serial dependencies
\cite{SekiEtAl1991,Kallmeyer2010}.

This paper fixes, as part of the learnable class, a homomorphism
$h:\Sigma^*\to M$ into a finite monoid and compares observed tuples only
when corresponding components have equal $h$-values.  The compositional type
guard is fixed before the positive text is seen, separating the choice of a
finite bias from learning under that bias~\cite{CosteTyping2004}.
Prefix--suffix and DFA transition morphisms provide two useful families.
Fan-out is part of the same bias because it determines which tuple arities are
tested for substitution.  A fixed bound on the unknown codomain size can be
absorbed into one product typing, whereas without such a bound the union
contains all regular languages and is not TxtEx-identifiable.  Thus the typing
supplies bounded finite-state information while the MCFG carries the
unbounded dependency.

Other learning biases make different choices.  Takada uses regular control
languages over derivations~\cite{Takada1988,Takada1995}, while local
substitutability restricts the evidence used to compare factors
\cite{CosteGaretNicolas2012,CosteNicolas2019}.  Positive-data learning has
also been studied for other mildly context-sensitive formalisms
\cite{OatesEtAl2006}.  MCFG learning with richer information includes
membership-query and dual context-set approaches
\cite{Yoshinaka2010,ClarkYoshinaka2014,KanazawaYoshinaka2023}.  The present
paper instead remains positive-data-only, uses observed tuple values as primal
objects, and lets a fixed finite typing guard substitutions.

\paragraph{Contributions.}
The paper has four main contributions.
\begin{enumerate}[label=\textup{(\arabic*)},leftmargin=2.4em]
\item \emph{A finite-typed form of Yoshinaka's positive-data learner.}
Yoshinaka's finite-sample grammar is recovered exactly under trivial typing
(up to notation and the separate treatment of $\eps$).  For a general fixed
finite typing, the reconstruction changes only by deleting substitution links
between tuples of unequal componentwise type.  A yield-typed refinement of a
target grammar proves that the guarded reconstruction is still complete after
a finite positive sample.  No rule-rank cap is needed for this qualitative
identification result; for each fixed cap $r_{\max}$, the corresponding
learner has polynomial-time updates in the accumulated positive sample.

\item \emph{A grammar-size polynomial at rule rank one.}
For arbitrary fixed $h$, rule-rank-one targets have characteristic data of
total length polynomial in grammar size alone.  The proof uses unary
derivation chains and applies to discontinuous tuples of arbitrary fixed
fan-out; no thickness parameter is needed.

\item \emph{A thickness-sensitive higher-rank transfer with an intrinsic-rank witness.}
At arbitrary fixed rank, typings that factor through a fixed prefix--suffix
summary admit a polynomial characteristic-data bound for reduced good
presentations in grammar size and rule thickness.  The proof simultaneously
preserves bounded boundary information in several tuple components.  The
rule-rank obstruction used below is inherited from the known non-branching
lower bound for $\mathrm{MIX}_2$; the contribution here is a single
$(1,1)$-typed non-context-free target combining that obstruction with minimum
fan-out two and a separation from Yoshinaka's trivial-typing class.  Its
minimum rule rank is two among \emph{all} finite-fan-out MCFG presentations,
so the higher-rank theorem applies beyond the rank-one regime at every finite
fan-out.
Full yield-type refinement can increase rule thickness exponentially, and
Appendix~\ref{app:normalization-details} shows the same phenomenon for
empty-component elimination.

\item \emph{Multiple-context expressiveness and fan-out sensitivity.}
For every $m\ge2$, a rank-one language of minimum fan-out $m$ belongs to the
$(1,1)$-typed class while lying outside Yoshinaka's class and the regular-filter
fragment recognized by that same typing.  A separate fan-out-two example lies
outside Yoshinaka's class followed by \emph{any} target-dependent regular
filter.  Increasing the tested tuple arity can already destroy substitutability
in Yoshinaka's untyped baseline; our stronger fan-out-two witness shows that
this loss need not be repairable by \emph{any} finite typing.
\end{enumerate}

\paragraph{Relation to prior work and scope.}
Yoshinaka supplies the semantic baseline, the observed-tuple grammar with
Type~I--III rules, the conservative update policy, rule-witness completeness,
and qualitative identification without a fixed rank bound
\cite[Sections~4.1--4.3, Corollary~2]{Yoshinaka2011}.  Accordingly,
the novelty here is the finite type guard and its MCFG-specific consequences
rather than a new observed-tuple grammar architecture.
Proposition~\ref{prop:yoshinaka-type-guard} gives the exact correspondence:
(Const) and (Comp) are Type~I, (Link) is Type~II with the type guard, and
(Start) is Type~III, apart from the separate $\eps$ convention.

Finite-monoid yield typing has a fan-out-one antecedent in the revised
companion CFG manuscript~\cite{kuriyamaCFG}, but all definitions and proofs
needed here are self-contained.  The MCFG-specific questions are componentwise
typing of tuples, rule-rank/thickness dependence of characteristic data, and
interaction with fan-out and regular filtering.  The non-branching lower bound
used for the intrinsic-rank witness is prior work; the contribution is its
boundary-typed wrapper, typed membership and Yoshinaka separation, and its
connection to the quantitative theorem.

Throughout the polynomial statements, fan-out, rank cap, and finite typing are
fixed.  Characteristic data for the set-driven reconstruction concern
language-level exactness; syntactic stabilization is supplied by the
conservative sequential learner.  The refinement obstructions concern the
displayed proof routes rather than lower bounds on all learners or
characteristic samples.

Sections~\ref{sec:prelim}--\ref{sec:char-data} develop the model,
reconstruction, convergence, and quantitative bounds; Section~\ref{sec:comparison}
gives the comparison with Yoshinaka's classes and the separation results.

\section{Preliminaries}
\label{sec:prelim}

Throughout, $\Sigma$ denotes a finite terminal alphabet, $\eps$ denotes the
empty word, $\Sigma^*$ is the set of all finite words over $\Sigma$, and
$\Sigma^+:=\Sigma^*\setminus\{\eps\}$.

\subsection{Positive-Data Learning}

We use Gold's explanatory learning model and restate the needed conventions
so that the MCFG development is self-contained.

\begin{definition}[Text and TxtEx identification]
We use Gold's explanatory identification criterion~\cite{Gold1967}.  To handle the empty language, fix a pause symbol
\(\pausesym\notin\Sigma\); thus
\(\pausesym,\pausesym,\ldots\) is a text for the empty language.
A \emph{text} for \(L\subseteq\Sigma^*\) is an infinite sequence over
\(L\cup\{\pausesym\}\) in which every element of \(L\) appears at least once.
Pause symbols may occur arbitrarily, and \(\pausesym\) is not an end-of-text marker.

The \emph{content} of a finite prefix is the finite set
\(K\subseteq L\) of positive examples from \(\Sigma^*\) that occur in that
prefix.  Thus neither the positions nor the number of pause symbols, nor repeated
occurrences of the same positive example, affect the content.
A learner is a computable function mapping each finite prefix to a hypothesis
grammar.

A learner \emph{TxtEx-identifies} \(L\) from positive data if, for every text
for \(L\), there are \(n_0\) and a single hypothesis grammar \(H\) such that
all stages from \(n_0\) onward output exactly \(H\), with \(L(H)=L\).  It
TxtEx-identifies a class when it does so for every language in the class.  A
sequential learner is \emph{conservative} if it keeps its current hypothesis
whenever the newly observed positive datum is already in that hypothesis
language, changing it only when the datum is not generated.
\end{definition}

For later use, fix a distinguished hypothesis code $H_0$ with
$L(H_0)=\emptyset$.

\begin{definition}[Set-driven reconstruction operator and characteristic sample]
\label{def:characteristic-sample}
A \emph{set-driven reconstruction operator} \(\mathcal B\) maps each finite
positive sample \(K\) to a grammar \(\mathcal B(K)\), depending only on the
set \(K\), not on presentation order or multiplicity.  It is
\emph{sample-consistent} if \(K\subseteq L(\mathcal B(K))\) for every
finite \(K\).

Following the characteristic-set viewpoint of de~la~Higuera
\cite{deLaHiguera1997}, for a target language \(L\), a finite set
\(K^\star\subseteq L\) is a \emph{characteristic sample for \(\mathcal B\)}
if, for every finite set
\(K\), \(K^\star\subseteq K\subseteq L\) implies
\(L(\mathcal B(K))=L\).  The case \(K^\star=\emptyset\) is allowed.  We
use \emph{characteristic data} as a neutral shorthand for the positive content
of such a sample; whenever a polynomial bound is claimed, we state explicitly
whether it concerns cardinality or total positive length.

This finite-sample notion is distinct from TxtEx learning: \(\mathcal B(K)\)
need not stabilize syntactically as the accumulated sample grows.  The
sequential learner will be introduced separately after exact reconstruction is
proved.
\end{definition}

\begin{definition}[Positive sample size]
For finite \(K\subseteq\Sigma^*\), define
\(\|K\|_+:=\sum_{w\in K}(|w|+1)\).
For a tuple \(\tuple u=(u_1,\ldots,u_d)\), write
\(\|\tuple u\|_1:=\sum_{i=1}^d |u_i|\)
for its total component length.
\end{definition}


\subsection{Finite-Monoid Typing}

\begin{definition}[Finite typing representation]
\label{def-monoid-morphism}
A finite-monoid homomorphism \(h\colon\Sigma^*\to M\) is represented by the
multiplication table of \(M\) together with the letter images \(h(a)\),
\(a\in\Sigma\).  Hence \(h(w)\) is computable by a left-to-right product in
linear time in \(|w|\).  For a tuple \(\tuple w=(w_1,\ldots,w_d)\), write
\(h^{(d)}(\tuple w):=(h(w_1),\ldots,h(w_d))\in M^d\).
\end{definition}

\begin{definition}[Refinement of finite typings]
\label{def:h-refinement}
Let \(h\colon\Sigma^*\to M\) and \(h'\colon\Sigma^*\to M'\) be finite
homomorphisms.  We say that \(h'\) \emph{refines} \(h\), and write
\(h\preceq h'\), when
\[
  \ker h'\subseteq\ker h.
\]
Equivalently, the value of \(h\) factors through the image of \(h'\): there is
a unique monoid homomorphism
\(\upsilon\colon\operatorname{im}h'\to\operatorname{im}h\) satisfying
\(h=\upsilon\circ h'\), after restricting the codomains to their images.
\end{definition}

For two typings $h$ and $g$, their product typing $(h,g)$ refines both and
satisfies $\ker(h,g)=\ker h\cap\ker g$.

Finite-monoid homomorphisms serve here as finite-state typings.  Let
$h_{\mathrm{triv}}:\Sigma^*\to\{1\}$ be the trivial typing.  A DFA
transition morphism for a regular superset of the target is another example
\cite{Pin1986}; only the induced finite summary is used, not a membership
oracle.  We call $h^{(d)}(\tuple w)$ the \emph{yield type}.  The
prefix--suffix morphism below supplies the boundary typing used later.

\begin{definition}[Prefix--suffix typing]
\label{def:prefix-suffix-typing}
For $k,\ell\ge0$, put $m_0:=\max\{1,k+\ell\}$.  For a word $w$, let
$\operatorname{pref}_k(w)$ and $\operatorname{suf}_\ell(w)$ denote its prefix
and suffix of the indicated lengths whenever $|w|\ge m_0$, with the length-zero
windows equal to $\eps$.  Define
\[
 h_\Sigma^{k,\ell}(w)=
 \begin{cases}
   w, & |w|<m_0,\\
   \langle\operatorname{pref}_k(w),\operatorname{suf}_\ell(w)\rangle,
      & |w|\ge m_0,
 \end{cases}
\]
using disjoint tags for the short-word and prefix--suffix elements.  Put
$M_{k,\ell}:=h_\Sigma^{k,\ell}(\Sigma^*)$.  For short words $x,y$ and long
elements $\langle u,v\rangle$, define
\[
\begin{aligned}
 x\cdot y&:=h_\Sigma^{k,\ell}(xy), &
 x\cdot\langle u,v\rangle&:=\langle\operatorname{pref}_k(xu),v\rangle,\\
 \langle u,v\rangle\cdot y&:=\langle u,\operatorname{suf}_\ell(vy)\rangle, &
 \langle u,v\rangle\cdot\langle u',v'\rangle&:=\langle u,v'\rangle.
\end{aligned}
\]
\end{definition}

With this multiplication,
$h_\Sigma^{k,\ell}(xy)=h_\Sigma^{k,\ell}(x)\cdot h_\Sigma^{k,\ell}(y)$.
Since $h_\Sigma^{k,\ell}$ is surjective, associativity is inherited from
concatenation and $h_\Sigma^{k,\ell}(\eps)$ is the identity.  Thus
$h_\Sigma^{k,\ell}:\Sigma^*\to M_{k,\ell}$ is a finite-monoid homomorphism.

A finite typing $h$ is called \emph{$(k,\ell)$-boundary-determined} when
$h\preceq h_\Sigma^{k,\ell}$.  Equivalently, equality of prefix--suffix types implies equality of $h$-types;
a factor map exists on the image of the prefix--suffix typing.


\subsection{Multiple Context-Free Grammars}

We use the standard tuple-generating definition of multiple context-free
grammars \cite{SekiEtAl1991}; see also \cite{Kallmeyer2010} for a general
overview of MCFGs and related mildly context-sensitive formalisms.
The general-rank notation below is used throughout the reconstruction and proofs.

\begin{definition}[Standard MCFG presentation]
\label{def:standard-mcfg}
A \emph{multiple context-free grammar} is a tuple
\(G=(V,\Sigma,P,S,\fan)\), where \(V\) is a finite set of nonterminals,
\(\Sigma\) is a finite terminal alphabet, \(S\in V\) has \(\fan(S)=1\), and
\(\fan(A)\ge1\) is the fan-out of each \(A\in V\).  A rule \(R\) of rank \(r\ge0\) has the form
\[
  R:\quad A\to\rho(B_1,\ldots,B_r),
  \qquad
  \rho=(\alpha_1,\ldots,\alpha_{d_0}),
\]
where \(\rho\) is the output template, \(d_0=\fan(A)\), \(d_j=\fan(B_j)\), and, with
\[
  \Xi_j:=\{\xi_j^1,\ldots,\xi_j^{d_j}\},
  \qquad
  \Gamma_\rho:=\Sigma\uplus\Xi_1\uplus\cdots\uplus\Xi_r,
\]
each \(\alpha_i\in\Gamma_\rho^*\) and every variable in
\(\Xi_1\uplus\cdots\uplus\Xi_r\) occurs at most once in the complete tuple
\((\alpha_1,\ldots,\alpha_{d_0})\).  We call this at-most-once condition
\emph{linearity}.  The rule is \emph{nondeleting} if every such variable occurs
exactly once.  Thus nondeleting does not mean that derived tuple components are
nonempty; empty strings remain allowed.

For child tuples
\(\tuple u_j=(u_{j,1},\ldots,u_{j,d_j})\in(\Sigma^*)^{d_j}\), simultaneous substitution of
\(u_{j,\kappa}\) for \(\xi_j^{\kappa}\) defines
\(\rho(\tuple u_1,\ldots,\tuple u_r)\in(\Sigma^*)^{d_0}.\)
When \(r=0\), the rule is a constant rule and derives the tuple
\((\alpha_1,\ldots,\alpha_{d_0})\).  The tuple languages \(L_A(G)\) form the least family, ordered componentwise by
inclusion, that is closed under all rule applications, and
\(L(G):=\{w\in\Sigma^*:(w)\in L_S(G)\}.\)
Equivalently, a finite derivation tree rooted at $A$ is evaluated bottom-up by
the rule templates.  We write $A\Rightarrow_G^*\tuple u$ when such a tree
yields $\tuple u\in(\Sigma^*)^{\fan(A)}$; this is a \emph{terminal
derivation}, and $\tuple u$ is its \emph{terminal tuple}.  A terminal
derivation $S\Rightarrow_G^*(w)$ from the start symbol is called
\emph{successful}.
An \(f\)-MCFG has \(\fan(A)\le f\) for every nonterminal, and an \(f\)-MCFL is
a language generated by an \(f\)-MCFG.  A grammar has \emph{rule-rank cap} \(r_{\max}\) if every rule has rank at most \(r_{\max}\).
\end{definition}

At fan-out one, the formalism includes ordinary context-free grammars.  A
\emph{linear context-free grammar} is a CFG in which every production has at
most one nonterminal on its right-hand side; a \emph{linear context-free
language} is generated by such a grammar.  This is distinct from the
``linearity'' condition on MCFG rule templates above, which means that each
rule variable occurs at most once.

Fix a standard explicit symbol-level encoding of finite MCFG presentations and
write $\|G\|$ for its length.  The encoding counts the nonterminal and rule
records and every terminal or variable occurrence in each rule template.  All
polynomial grammar-size bounds below refer to this encoding.

For integers \(f\ge1\) and \(r\ge0\), write
\[
\begin{aligned}
  \mathsf{MCFL}(f,r):=\{L(G):{}&G\text{ is an MCFG of fan-out at most }f,\\
                   &\text{with rule rank at most }r\}.
\end{aligned}
\]
Thus \(f\text{-}\mathrm{MCFL}=\bigcup_{r<\infty}\mathsf{MCFL}(f,r)\), because every
finite MCFG presentation has finite rule rank.
When minimum rule rank at fan-out \(f\) is mentioned below, it means the
language-level minimum over all fan-out-\(f\) presentations, not the rank of a
particular supplied grammar.


We use $f$ throughout for fan-out/dimension and $r$ for a generic rule rank;
$r_{\max}$ denotes a prescribed rank cap.  Rule-template variables are always
written $\xi_j^\kappa$, so that $x,y,z$ (with or without arrows) remain free for
words and observed tuples.  Terminals in examples, including digits and
$\mathtt\#$, are typeset in \texttt{typewriter} font.


When two nonterminals have the same fan-out \(d\), we use
\(A\to B\)
as shorthand for the unary identity-template rule
\(A\to(\xi_1^1,\ldots,\xi_1^d)(B)\).  This denotes only the identity
correspondence between components: at fan-out two, for example, it means
\(A\to(\xi_1^1,\xi_1^2)(B)\), not a swapping rule such as
\(A\to(\xi_1^2,\xi_1^1)(B)\).
Likewise, \(A\to\tuple c\) denotes a rank-zero
constant rule.  These abbreviations are used for typed start links and for Rule~(Link) of the
reconstruction below.

The standard formalism allows rank zero, nonidentity unary rules, arbitrary
finite rule rank, and deleting rules.  A nondeleting rule
\(R:A\to\rho(B_1,\ldots,B_r)\), with
\(\rho=(\alpha_1,\ldots,\alpha_{d_0})\), is
\emph{nonpermuting} if, for each child \(B_j\), the variables
\(\xi_j^1,\ldots,\xi_j^{\fan(B_j)}\) occur in this order when the parent components
are scanned from left to right.  Following Yoshinaka
\cite[Section~2.3]{Yoshinaka2011}, a rule is \emph{$\lambda$-free} if no output
template $\alpha_i$ is empty, and a nonpermuting rule is \emph{nonmerging} if
no $\alpha_i$ contains adjacent variables $\xi_j^{\kappa} \xi_j^{\kappa+1}$ of the same child.
A presentation is \emph{good} if all its rules are $\lambda$-free,
nondeleting, nonpermuting, and nonmerging.  We write $\eps$ for the empty word
throughout; ``$\lambda$-free'' is retained only as Yoshinaka's standard name
for this rule condition.  We use good presentations in the completeness proof.  The nonmerging condition ensures that, after the siblings
of a child are fixed, consecutive components of that child remain separated
by nonempty terminal material or by a nonempty sibling yield.  This matches
Yoshinaka's multidimensional-context convention.  In particular,
$\lambda$-freeness implies by induction that every component of every terminal tuple derived from a
nonterminal is nonempty.

The normalization used below uses only the nondeleting/information-lossless
step (f3) in the proof of Seki et al.'s Lemma~2.2, in the rank-preserving form
recorded by Kanazawa, together with component permutation, explicit
empty-component removal, and Yoshinaka's nonmerging conversion
\cite[proof of Lemma~2.2, item~(f3)]{SekiEtAl1991}\cite[Section~2.1]{Kanazawa2019Ogden}\cite[Section~2.3, Lemma~1]{Yoshinaka2011}.  Yoshinaka also notes later in
Section~2.3, after Example~2, that for fixed dimension and rule rank these
standard conversions yield
an equivalent good grammar of size $O(\|G\|)$~\cite[Section~2.3]{Yoshinaka2011}.

Fan-out and rule rank are distinct grammar resources.  In particular, there is
no blanket binarization theorem that preserves an arbitrary fixed fan-out for
general MCFG/LCFRS presentations: reducing rule rank may require increased
fan-out \cite{RambowSatta1999,GomezRodriguezEtAl2009}.
Accordingly the qualitative reconstruction below is defined without imposing a
binary normal form or a known rank cap.


\paragraph{Reduced grammars.}
The nonterminal dependency graph of an MCFG has an edge $A\to B$ whenever
$B$ occurs as a child in a rule with left-hand side $A$.
A nonterminal is \emph{reachable} if it is reachable from the start symbol in
this graph, and \emph{productive} if its tuple language is nonempty.
The grammar is \emph{reduced} if every nonterminal is both reachable and
productive.


\subsection{Normalization}

\begin{proposition}[Good normalization with rank preservation]
\label{prop:good-normalization}
Let $G_{\mathrm{in}}$ be a finite standard $f$-MCFG of maximum rule rank at
most $r_{\max}$.  There is a finite good $f$-MCFG $G_{\mathrm{good}}$ of
maximum rule rank at most $r_{\max}$ such that
\[
  L(G_{\mathrm{good}})=L(G_{\mathrm{in}})\setminus\{\eps\}.
\]
For fixed $f$ and $r_{\max}$,
$\|G_{\mathrm{good}}\|=O_{f,r_{\max}}(\|G_{\mathrm{in}}\|)$.
If $\eps\notin L(G_{\mathrm{in}})$, the two grammars are equivalent.
\end{proposition}

\begin{proof}
First delete unreachable or unproductive nonterminals, then apply the
nondeleting/information-lossless step (f3) from the proof of Seki et al.'s
Lemma~2.2, in the rank-preserving form recorded by
Kanazawa~\cite[proof of Lemma~2.2, item~(f3)]{SekiEtAl1991}\cite[Section~2.1]{Kanazawa2019Ogden}.
Decorating nonterminals by component permutations makes the grammar
nonpermuting.  Next decorate by nonempty-component patterns realizable by terminal
derivations, as made explicit in Appendix~\ref{app:normalization-details}, and
project away empty components; this removes $\eps$ at the start and makes every
retained output template $\lambda$-free.  Finally apply Yoshinaka's nonmerging
conversion~\cite[Section~2.3, Lemma~1]{Yoshinaka2011}.
None of these steps increases fan-out or rule rank.  With $f,r_{\max}$ fixed,
component-pattern choices contribute at most $2^f$ nonterminal decorations and
$2^{f r_{\max}}$ child-pattern choices per rule, while permutation decoration
contributes at most $f!$ copies.  Yoshinaka's conversion has the same
fixed-parameter linear-size behavior.  Hence the total encoding size remains
$O_{f,r_{\max}}(\|G_{\mathrm{in}}\|)$.  The derivation correspondence and the
empty-component construction are recorded in
Appendix~\ref{app:normalization-details}.
\end{proof}

This proposition controls size, fan-out, and rule rank, but it does not
identify the rule thickness of an arbitrary presentation with the thickness
after normalization.  Thickness-sensitive statements are therefore first
proved for good presentations.  For reduced $\lambda$-free input
presentations, including presentations with deleting rules,
Proposition~\ref{prop:structural-good-thickness} gives a polynomial
grammar-size-and-thickness transfer to the structural good form used below.

\subsection{Tuple Contexts and Substitutability}

\paragraph{Sentence contexts.}
A \emph{sentence context of arity $d$} is a word
$E=v_0\blank_1v_1\cdots v_{d-1}\blank_dv_d$ in which each named hole occurs
exactly once, $v_0,v_d\in\Sigma^*$, and every internal separator
$v_1,\ldots,v_{d-1}$ belongs to $\Sigma^+$.  Let $\mathcal E_d$ be the set
of such contexts.  This is Yoshinaka's multidimensional-context convention
\cite[Section~3.1]{Yoshinaka2011}.
For $\tuple w=(w_1,\ldots,w_d)$, write $E[\tuple w]$ for the result of filling
$\blank_i$ by $w_i$, and put
$|E|_{\mathrm{fix}}:=\sum_{i=0}^d|v_i|$.  For brevity, we say that a
sentence context $E$ \emph{accepts} $\tuple w$ in a language $L$ when
$E[\tuple w]\in L$.

\begin{definition}[Tuple occurrence in a sample word]
\label{def:tuple-occurrence}
Let $K\subseteq\Sigma^*$ and $w=a_1\cdots a_n\in K$.  Number the cut
positions of $w$ as
$\cutpos{0}a_1\cutpos{1}\cdots a_n\cutpos{n}$.
For $0\le i<i'\le n$, write
$w[i:i']:=a_{i+1}\cdots a_{i'}$.

An \emph{arity-$d$ tuple occurrence} in $w$ is a sequence of nonempty,
left-to-right intervals
\[
  0\le i_1<i_1'<i_2<i_2'<\cdots<i_d<i_d'\le n.
\]
It determines the tuple
$\tuple x=(w[i_1:i_1'],\ldots,w[i_d:i_d'])\in(\Sigma^+)^d$ and the
context $E\in\mathcal E_d$ obtained by replacing those intervals by
$\blank_1,\ldots,\blank_d$.  We identify the occurrence with $(E,\tuple x)$.
For example, the two intervals spelling $\mathtt a$ and $\mathtt c$ in $\mathtt{abc}$ give the tuple
$(\mathtt a,\mathtt c)$ and the context $\blank_1\mathtt b\blank_2$.
\end{definition}

For $L\subseteq\Sigma^*$ and $\tuple x\in(\Sigma^*)^d$, define its
\emph{$d$-tuple distribution}
\[
  \D_L^{(d)}(\tuple x):=\{E\in\mathcal E_d:E[\tuple x]\in L\}.
\]
When $d$ is clear, the superscript is omitted.

\begin{definition}[$(f,h)$-tuple substitutability]
\label{def:tuple-substitutable}
Let $f\ge1$ and let $h:\Sigma^*\to M$ be a finite monoid
homomorphism.  A language $L\subseteq\Sigma^*$ is
\emph{$(f,h)$-tuple-substitutable} if, for every $1\le d\le f$ and all
$\tuple x,\tuple y\in(\Sigma^+)^d$,
\[
 h^{(d)}(\tuple x)=h^{(d)}(\tuple y),\qquad
 \D_L^{(d)}(\tuple x)\cap\D_L^{(d)}(\tuple y)\ne\emptyset
\]
imply $\D_L^{(d)}(\tuple x)=\D_L^{(d)}(\tuple y)$.
Thus empty strings are not compared as tuple components; when $\eps$ belongs
to the target, it is handled separately at the start symbol.  This is Yoshinaka's nonempty-component convention for multidimensional
substitutability~\cite[Section~3.1 and footnote~3]{Yoshinaka2011}; see also Clark~\cite{Clark2013}.  The restriction is
substantive already at fan-out one: if $\eps$ itself were allowed among the
compared fragments, the ordinarily substitutable language $\mathtt a^+$ would fail
the condition because $\mathtt a$ and $\eps$ share a context but do not have the same
distribution.
\end{definition}

At fan-out one, every sentence context has the form
$u\blank_1v$, so $\D_L^{(1)}(x)$ is canonically the ordinary two-sided
distribution $\D_L(x)=\{(u,v):uxv\in L\}$.  Hence Definition~\ref{def:tuple-substitutable}
specializes, without any external convention, to the ordinary two-sided
fixed-$h$ substitutability condition on nonempty factors.

\paragraph{Typing--distribution equivalence.}
For $\tuple x,\tuple y\in(\Sigma^+)^d$, write
$\tuple x\equiv_{L,h}^d\tuple y$ if and only if
\[
 h^{(d)}(\tuple x)=h^{(d)}(\tuple y)
 \qquad\text{and}\qquad
 \D_L^{(d)}(\tuple x)=\D_L^{(d)}(\tuple y).
\]
For fixed $L,h,d$ this is an equivalence relation; the superscript is omitted
when the arity is clear.

\begin{definition}[The target class]
\label{def:class-cmcf}
Fix $f\ge1$ and a finite monoid homomorphism $h$.  The class
$\Cmcf{f}{h}$ consists of all $f$-MCFLs that are
$(f,h)$-tuple-substitutable.
\end{definition}

If $h\preceq h'$, then
$\Cmcf{f}{h}\subseteq\Cmcf{f}{h'}$: equality of the finer $h'$-types implies
equality of the coarser $h$-types, so the substitution conclusion available
under $h$ applies.  At fan-out one this is the context-free part of the
two-sided fixed-$h$ condition just described.

For every $f\ge1$,
\[
  \Cmcf{f+1}{h}\cap f\text{-}\mathrm{MCFL}
  \subseteq \Cmcf{f}{h},
\]
because the fan-out-$f$ condition tests only the arities
$1,\ldots,f$.  The possible loss of substitutability when the tested arity is
increased is already present for trivial typing: Yoshinaka notes that
$\mathsf{SL}(p,r)$ and $\mathsf{SL}(p+1,r)$ are generally incomparable and
gives
$L_{\mathrm{reverse}}=\{w\#w^R\}$ as a one-dimensional substitutable language
that is not two-dimensional substitutable
\cite[Section~4.1 and Example~5]{Yoshinaka2011}.  What finite
typing adds here is the stronger phenomenon of
Proposition~\ref{prop:strong-fanout-separation}: a target can be learnable at
fan-out one under a finite typing yet fail the fan-out-two condition for
\emph{every} finite typing.  Thus $f$ is part of the learning bias, not merely
an upper bound on representational power.  No monotonicity of the learned
classes in $f$ is claimed.

For every $\mathsf{Acc}\subseteq M$, the language
$h^{-1}(\mathsf{Acc})$ belongs to $\Cmcf{f}{h}$.  Indeed, it is regular, and equal componentwise $h$-types give the
same total $h$-value in every sentence context, hence equal tuple
distributions.

\begin{proposition}[Intersection and abelian filtering]
\label{prop:abelian-coset-filter}
Fix $f$ and $h$.
\begin{enumerate}[label=(\roman*),leftmargin=*]
\item If $L_1$ and $L_2$ are $(f,h)$-tuple-substitutable, then
$L_1\cap L_2$ is $(f,h)$-tuple-substitutable.
\item Let $L$ be $(f,h)$-tuple-substitutable, let $\mathbb A$ be an abelian
group, let $\varphi:\Sigma^*\to\mathbb A$ be a monoid homomorphism, let
$\mathcal H\le\mathbb A$, and let $\gamma\in\mathbb A$.  Then
$L\cap\varphi^{-1}(\gamma+\mathcal H)$ is $(f,h)$-tuple-substitutable.
\end{enumerate}
Whenever the resulting intersection is an $f$-MCFL, it therefore belongs to
$\Cmcf{f}{h}$.
\end{proposition}

\begin{proof}
For (i), equal-type tuples $\tuple x,\tuple y$ with some context $E$ satisfying
$E[\tuple x],E[\tuple y]\in L_1\cap L_2$ have equal tuple distributions in each $L_i$, hence in their intersection.
For (ii), first observe that $\varphi^{-1}(\gamma+\mathcal H)$ is
$(f,h_{\mathrm{triv}})$-tuple-substitutable.  Indeed, if a common context
accepts tuples $\tuple x,\tuple y$, then
$\sum_i(\varphi(y_i)-\varphi(x_i))\in\mathcal H$.  By commutativity, for every other sentence
context $E'$ the difference
$\varphi(E'[\tuple y])-\varphi(E'[\tuple x])$ is the same element of
$\mathcal H$, so coset membership is preserved in both directions.  Trivial
typing is coarser than $h$, hence the same coset language is
$(f,h)$-tuple-substitutable.  Apply (i).
\end{proof}

As an immediate consequence, for any $\mathsf{Acc}\subseteq M$, abelian group
$\mathbb A$, homomorphism $\varphi:\Sigma^*\to\mathbb A$, subgroup
$\mathcal H\le\mathbb A$, and $\gamma\in\mathbb A$, every $f$-MCFL of the
form
\[
 h^{-1}(\mathsf{Acc})\cap\varphi^{-1}(\gamma+\mathcal H)
\]
belongs to $\Cmcf{f}{h}$.  Thus one typing can support a family of targets
while the abelian constraint varies; these factorizations are not supplied to
the learner.

\section{Exact Reconstruction from Positive Data}
\label{sec:learner}

We define the finite-sample reconstruction operator, prove soundness, and then
prove completeness using a normalized target grammar refined by yield type.
The target grammar is used only to select a finite characteristic sample.
Since every observed tuple component is nonempty, the rank of an enumerated
composition witness is bounded by the length of the sample word containing it.

\subsection{Set-Driven Reconstruction and Sample-Bounded Rank}


Following Definition~\ref{def:tuple-occurrence}, enumerate all tuple occurrences
of arity at most $f$ in the sample $K$.  The reconstruction uses one nonterminal
$[\tuple x]$ for every observed tuple value $\tuple x\in(\Sigma^+)^d$, together
with a fresh start symbol $\widehat S$.  It is not told which occurrences
correspond to target nonterminals or rules; all admissible segmentations are
enumerated.

\begin{definition}[Composition witness]
\label{def:composition-witness}
Let $r\ge1$.  A \emph{composition witness of rank $r$} in $K$ consists of an
observed parent occurrence $(E_{\mathrm{par}},\tuple z)$ of arity $e\le f$ together with a
segmentation of the parent components into terminal gaps and nonempty labelled
intervals.  The labels are variables $\xi_j^{\kappa}$ with
$1\le j\le r$ and $1\le\kappa\le d_j\le f$; every variable occurs exactly once,
and for each fixed $j$ the variables
$\xi_j^1,\ldots,\xi_j^{d_j}$ occur in this order when the parent components are
scanned from left to right.  No output component may contain adjacent variables
$\xi_j^{\kappa}\xi_j^{\kappa+1}$ of the same child.  The labelled intervals
induce child tuples $\tuple u_j=(u_{j,1},\ldots,u_{j,d_j})$ and determine a
unique linear template $\rho$ satisfying
\[
  \tuple z=\rho(\tuple u_1,\ldots,\tuple u_r).
\]
\end{definition}

\begin{lemma}[Materialization of child tuples]
\label{lem:composition-witness-materialization}
Every child tuple induced by a composition witness is itself an observed tuple
occurrence, and the induced template $\rho$ is good.
\end{lemma}

\begin{proof}
Because the parent is an observed tuple occurrence, the components of each
$\tuple u_j$ occur in their original left-to-right order in the same sample
word.  If two consecutive components of the same child lie in one parent
component, the definition of a composition witness excludes their direct
adjacency.  They are therefore separated either by nonempty terminal material
or by one or more nonempty labelled intervals belonging to other children.
If they lie in different parent components, they are separated by the
nonempty internal material of the parent sentence context.  Hence the
components of each $\tuple u_j$ satisfy the separation condition for an
observed tuple occurrence.

Every variable occurs exactly once, the components of each child occur in
order, and adjacent consecutive components of one child are excluded.
Therefore $\rho$ is linear, nondeleting, nonpermuting, and nonmerging.
Finally every parent component is nonempty, so every output component of
$\rho$ is nonempty.  Thus $\rho$ is good.
\end{proof}

\begin{lemma}[Enumeration and sample-bounded rank]
\label{lem:general-witness-enumeration}
Fix $f$ and a finite sample $K$, and put $n_K:=\max\{1,\|K\|_+\}$.  If a
rank-$r$ composition witness has parent tuple $\tuple z$ occurring in a sample
word $w$, then
\[
  r\le\|\tuple z\|_1\le |w|\le\|K\|_+.
\]
All composition witnesses are effectively enumerable.  Under a rule-rank cap
$r_{\max}$, those of rank at most $r_{\max}$ are enumerable in
$n_K^{O(f r_{\max})}$ time and space.
\end{lemma}

\begin{proof}
Every child tuple has at least one nonempty component.  Linearity and
nondeletion make the child components contribute disjointly to the parent, so
$\|\tuple z\|_1\ge\sum_j\|\tuple u_j\|_1\ge r$; the remaining inequalities
follow because the parent components are disjoint intervals of $w$.

An arity-$d$ tuple occurrence is determined by $2d$ cuts.  For a fixed rank
$r$, a candidate segmentation has at most $fr$ labelled child intervals, whose endpoints are chosen among the $|w|+1\le n_K+1$ cut positions of
the parent sample word.  For $r\le r_{\max}$, assignments to parent components and
terminal gaps contribute only a factor depending on $f$ and $r_{\max}$, and
the nonpermuting and nonmerging conditions are directly checkable.  Summing
over $r\le r_{\max}$ gives the stated bound.  Without a cap, the displayed
rank inequality restricts the effective enumeration to
$1\le r\le\|K\|_+$.
\end{proof}

\begin{definition}[Capped and uncapped reconstruction operators]
\label{def:reconstruction}
Fix $f$ and $h$.  For $r_{\max}\ge1$, let
$\mathcal B_{f,h}^{\le r_{\max}}(K)$ be the grammar
$\widehat G_{\le r_{\max}}(K)$ with a fresh start symbol $\widehat S$, one nonterminal $[\tuple x]$ for each nonempty tuple value $\tuple x$ observed in $K$,
and the following rules:
\begin{enumerate}[leftmargin=*]
\item[\textnormal{(Start)}] for every $w\in K\cap\Sigma^+$,
      $\widehat S\to[(w)]$, and if $\eps\in K$, also
      $\widehat S\to\eps$;
\item[\textnormal{(Const)}] for every observed tuple, the constant rule
      $[\tuple u]\to\tuple u$;
\item[\textnormal{(Comp)}] for every composition witness of rank
      $1\le r\le r_{\max}$ with parent tuple $\tuple z$, template $\rho$,
      and induced child tuples $\tuple u_1,\ldots,\tuple u_r$, the rule
      \[
        [\tuple z]\to\rho([\tuple u_1],\ldots,[\tuple u_r]);
      \]
\item[\textnormal{(Link)}] for observed nonempty tuples
      $\tuple x,\tuple y$ of the same arity $d$, the identity-template rule
      $[\tuple x]\to[\tuple y]$ whenever
      $h^{(d)}(\tuple x)=h^{(d)}(\tuple y)$ and some $E\in\mathcal E_d$
      satisfies $E[\tuple x],E[\tuple y]\in K$.
\end{enumerate}
If $K\ne\emptyset$, define
$\mathcal B_{f,h}(K):=\widehat G(K):=
\widehat G_{\le\|K\|_+}(K)$; for $K=\emptyset$, let
$\mathcal B_{f,h}(\emptyset)$ be the grammar with no productive start rule.
Fixed encoding orders make both reconstruction operators deterministic and
set-driven.
\end{definition}

\paragraph{Exact relation to Yoshinaka's finite-sample grammar.}
The finite typing modifies Yoshinaka's constructor only at substitution links.
\begin{proposition}[Typed guard on Yoshinaka's reconstruction]
\label{prop:yoshinaka-type-guard}
Fix a finite positive sample $K$, fan-out bound $f$, and rank cap
$r_{\max}\ge1$.  Under the trivial typing,
$\mathcal B_{f,h_{\mathrm{triv}}}^{\le r_{\max}}(K)$ agrees, up to notation,
with Yoshinaka's hypothesis grammar $G(K)$ for $\mathcal A(f,r_{\max})$
\cite[Section~4.1]{Yoshinaka2011}, except that the present construction handles
$\eps\in K$ by a separate direct start rule.  For an arbitrary finite typing
$h$, $\mathcal B_{f,h}^{\le r_{\max}}(K)$ is obtained from that grammar by
deleting exactly those Type~II substitution rules whose source and target
tuples have different componentwise $h$-types.  The same statement holds for
the uncapped reconstruction and Yoshinaka's $\mathcal A(f,*)$.
\end{proposition}

\begin{proof}
The observed tuple values are the same nonterminal indices in both
constructions.  Rule~(Start) is Yoshinaka's Type~III rule, and Rule~(Link) is
Type~II with the additional componentwise type test.

It remains to compare Type~I with Rules~(Const) and~(Comp).  A rank-zero
Type~I rule is exactly Rule~(Const).  For positive rank, suppose Yoshinaka has
\[
 [\tuple z]\to\rho([\tuple u_1],\ldots,[\tuple u_r])
 \qquad\text{with}\qquad
 \tuple z=\rho(\tuple u_1,\ldots,\tuple u_r),
\]
where all displayed tuple values are observed and $\rho$ is good.  Choose an
observed occurrence of $\tuple z$ in a sample word.  The displayed equality
locates every component of every $\tuple u_j$ as a labelled interval inside
that occurrence.  Linearity and nondeletion use each interval once,
nonpermutation puts the components of each child in order, and nonmerging
prevents direct adjacency of consecutive components of one child.  Hence these
intervals form a composition witness, so Rule~(Comp) contains the same rule.
Conversely every composition witness is a good algebraic decomposition among
observed tuple values and therefore is a Type~I rule.  This also matches
Yoshinaka's counting proof, where each Type~I rule is enumerated as a
factorization of a single sample word
\cite[Lemma~6]{Yoshinaka2011}.

Thus trivial typing gives the same Type~I--III rule set, apart from the direct
$\eps$ start convention.  For general $h$, Rules~(Start), (Const), and (Comp)
are unchanged, and only Type~II/Rule~(Link) is filtered by componentwise type.
Finally, in the uncapped construction every composition witness has rank at
most the length of the sample word containing it, so allowing the cap
$\|K\|_+$ introduces no Type~I rule beyond the sample-length cap used in
Yoshinaka's $\mathcal A(f,*)$.
\end{proof}

\begin{example}[Why the equal-type premise matters]
\label{ex:type-guard-purpose}
Let $L=\{\mathtt a^n\mathtt b^n\mathtt c^n:n\ge1\}$ and
$K=\{\mathtt{abc},\mathtt{aabbcc}\}\subseteq L$.  The observed factors
$(\mathtt b)$ and $(\mathtt{abbc})$ share the sample context
$\mathtt a\blank_1\mathtt c$.  With trivial typing, Rule~(Link) therefore
connects their hypothesis nonterminals.  The composition witness
$[\mathtt{aabbcc}]\to\mathtt{aa}[\mathtt b]\mathtt{bcc}$ can then use that
link to derive $\mathtt{aaabbcbcc}\notin L$.  Under $h_\Sigma^{1,1}$ the
link is blocked because $h_\Sigma^{1,1}(\mathtt b)=\mathtt b$ while
$h_\Sigma^{1,1}(\mathtt{abbc})=\langle\mathtt a,\mathtt c\rangle$.
This illustrates the local role of the equal-type premise; the strictness theorem in
Section~\ref{sec:comparison} later uses the same boundary mechanism at every
fan-out level.
\end{example}

The reconstruction is sample-consistent:
\[
 K\subseteq L\bigl(\mathcal B_{f,h}^{\le r_{\max}}(K)\bigr),
 \qquad
 K\subseteq L\bigl(\mathcal B_{f,h}(K)\bigr).
\]
Indeed, each nonempty $w\in K$ has the direct derivation
$\widehat S\to[(w)]\to(w)$ by Rules~(Start)--(Const), while
$\eps\in K$ gives the direct start rule $\widehat S\to\eps$.

These $\mathcal B$ operators are set-driven constructors, not sequential TxtEx
learners: recomputation on growing samples need not stabilize syntactically
(e.g. for successively longer $\mathtt a^n\mathtt c\mathtt b^n$).  The conservative learner below
rebuilds only when a new positive datum is rejected.

Both reconstruction operators produce finite standard MCFGs of fan-out at
most $f$ and are effectively constructible from $K,f,h$.  There are only
finitely many observed nonempty tuple occurrences.  Rule~(Link) is obtained by
grouping occurrences by concrete context and componentwise type, while
Lemma~\ref{lem:general-witness-enumeration} enumerates Rule~(Comp).  By
Lemma~\ref{lem:general-witness-enumeration}, every rule of $\mathcal B_{f,h}(K)$
has rank at most $\|K\|_+$ when $K\ne\emptyset$.

The concrete separation examples in Section~\ref{sec:comparison} also serve as
small instances of the tuple-nonterminal and composition mechanisms used here.

\subsection{Soundness}
\label{sec:soundness}

\begin{lemma}[Induced child context]
\label{lem:filling-identity}
Let $\rho$ be the template of a good rule, let
$E\in\mathcal E_e$ be a sentence context for its output tuple, and fix tuples of the arities required by the child nonterminals, with nonempty
components, for all children except child $j$.
Replacing every variable of child $j$ by its named hole and every other child
variable by the corresponding fixed component produces a context
$E_j\in\mathcal E_{d_j}$
such that
\[
 E_j[\tuple u_j]=E[\rho(\tuple u_1,\ldots,\tuple u_r)].
\]
If the sibling tuples are changed, the same construction gives the
corresponding induced context $E_j$ for that new sibling choice; the context
itself need not be identical.
\end{lemma}

\begin{proof}
Linearity and nondeletion make every variable of child $j$ occur exactly once,
and nonpermutingness puts $\xi_j^1,\ldots,\xi_j^{d_j}$ in that order.  Between two
consecutive variables of this child, nonmerging ensures either a terminal
symbol, a variable of another child, or a boundary between two output
components.  In the first case the separator is visibly nonempty; in the
second it becomes a nonempty sibling component; in the third the corresponding
internal separator of $E$ is nonempty.  Replacing the distinguished variables
by $\blank_1,\ldots,\blank_{d_j}$ therefore gives a context in
$\mathcal E_{d_j}$, and filling the holes is exactly composition with $\rho$
and then with $E$.
\end{proof}

\begin{lemma}[Composition preserves equivalence]
\label{lem:witnessed-composition}
Let $L\subseteq\Sigma^*$, let $\rho$ be the template of a good rank-$r$ rule, and let
$\tuple x_j,\tuple u_j$ have the arities required by its children and have
nonempty components.  If
$\tuple u_j\equiv_{L,h}\tuple x_j$ for every $j$, then
\[
 \rho(\tuple u_1,\ldots,\tuple u_r)
 \equiv_{L,h}
 \rho(\tuple x_1,\ldots,\tuple x_r).
\]
\end{lemma}

\begin{proof}
Homomorphic template evaluation gives equality of the two parent $h$-types,
and goodness of $\rho$ makes every parent component nonempty.  It remains to
compare distributions.  Let $E$ be an arbitrary sentence context for the
parent tuple.  If
$E[\rho(\tuple x_1,\ldots,\tuple x_r)]\in L$, replace the children one at
a time.  At step $j$, Lemma~\ref{lem:filling-identity} supplies the induced
sentence context for child $j$ under the current sibling choice.  Equality of
the child distributions transfers acceptance from $\tuple x_j$ to
$\tuple u_j$.  After all replacements,
$E[\rho(\tuple u_1,\ldots,\tuple u_r)]\in L$.  Reversing the replacements
gives the converse implication.  Since $E$ was arbitrary, the two parent
tuples have equal distributions.
\end{proof}

\begin{proposition}[Soundness of every capped hypothesis]
\label{prop:reconstruction-soundness}
Let $L$ be $(f,h)$-tuple-substitutable and $K\subseteq L$.  For every
$r_{\max}\ge1$, if $[\tuple x]\Rightarrow^*\tuple u$ in
$\widehat G_{\le r_{\max}}(K)$, then every component of $\tuple u$ is nonempty
and $\tuple u\equiv_{L,h}\tuple x$.  Consequently
\[
 L\bigl(\mathcal B_{f,h}^{\le r_{\max}}(K)\bigr)\subseteq L,
 \qquad L\bigl(\mathcal B_{f,h}(K)\bigr)\subseteq L.
\]
\end{proposition}

\begin{proof}
Induct on derivation height.  Constant rules are immediate.  For a Rule~(Link) instance $[\tuple x]\to[\tuple y]$, its defining common sample context and
equal $h$-type give $\tuple x\equiv_{L,h}\tuple y$; transitivity with the
induction hypothesis proves the claim.  For a composition-witness rule, the
induction hypotheses give nonempty child tuples equivalent to the witnessed
children.  Every parent component of the witnessed tuple was nonempty; under
the same nondeleting template it remains nonempty after replacing child
components by other nonempty strings.  Lemma~\ref{lem:witnessed-composition}
then gives the parent equivalence.

A nonempty start derivation begins with $\widehat S\to[(w_0)]$ for some
$w_0\in K\cap\Sigma^+$, so the empty unary context places the derived word in
$L$.  The direct rule $\widehat S\to\eps$ is present only when
$\eps\in K\subseteq L$.  The $K=\emptyset$ case is trivial, and the general
hypothesis is the capped case with $r_{\max}=\|K\|_+$ whenever $K\ne\emptyset$.
\end{proof}


\subsection{Yield-Typed Target Refinement and Finite Witnesses}
\label{sec:refinement}

\begin{definition}[Template yield type]
\label{def:output-map}
Let
\(R:A\to\rho(B_1,\ldots,B_r)\)
be a nondeleting rule with output template
\(\rho=(\alpha_1,\ldots,\alpha_{d_0})\), and put \(d_0=\fan(A)\) and
\(d_j:=\fan(B_j)\), and choose
\(\mathbf t_j\in M^{d_j}\) for each \(1\le j\le r\).

On the template alphabet
\(
  \Gamma_\rho
  :=\Sigma\uplus\biguplus_{j=1}^r
  \{\xi_j^1,\ldots,\xi_j^{d_j}\}
\),
define
\(\evalmap_{\mathbf t_1,\ldots,\mathbf t_r}\colon\Gamma_\rho\to M\)
by
\(\evalmap_{\mathbf t_1,\ldots,\mathbf t_r}(a):=h(a)\)
for \(a\in\Sigma\), and
\(\evalmap_{\mathbf t_1,\ldots,\mathbf t_r}(\xi_j^{\kappa}):=(\mathbf t_j)_\kappa\).

Let
\(
  \widehat{\evalmap}_{\mathbf t_1,\ldots,\mathbf t_r}
  \colon\Gamma_\rho^*\to M
\)
be its unique extension to a monoid homomorphism.
Define
\[
  \out_\rho(\mathbf t_1,\ldots,\mathbf t_r)
  :=
  \bigl(
    \widehat{\evalmap}_{\mathbf t_1,\ldots,\mathbf t_r}(\alpha_1),
    \ldots,
    \widehat{\evalmap}_{\mathbf t_1,\ldots,\mathbf t_r}(\alpha_{d_0})
  \bigr)
\in M^{d_0}.
\]
For rank zero, \(\out_\rho\) is simply
\((h(\alpha_1),\ldots,h(\alpha_{d_0}))\), the componentwise \(h\)-type of the
constant template.
\end{definition}

The refinement below is ordinary finite nonterminal splitting: compositionality
$h(uv)=h(u)h(v)$ makes each typed target nonterminal uniform in componentwise
$h$-type, which is exactly the guard needed for substitutability.


\begin{definition}[Yield-type refinement]
\label{def:output-refinement}
Given a finite nondeleting MCFG \(G=(V,\Sigma,P,S,\fan)\) and a finite
typing \(h\), the \emph{full yield-type refinement} \(G^h\) has a fresh
start symbol \(\widetilde S\) and typed nonterminals
\(A_{\mathbf t}\) for every \(A\in V\) and
\(\mathbf t\in M^{\fan(A)}\).
Set
\(\fan(\widetilde S):=1\) and
\(\fan(A_{\mathbf t}):=\fan(A)\).

For each \(\eta\in M\), add the start rule
\(\widetilde S\to S_{(\eta)}\).
For each rank-zero rule \(A\to\tuple c\), add
\(A_{h^{(\fan(A))}(\tuple c)}\to\tuple c\).

For every positive-rank rule
\[
R:A\to\rho(B_1,\ldots,B_r),
\qquad
\rho=(\alpha_1,\ldots,\alpha_{d_0}),
\]
and every choice $\mathbf t_j\in M^{\fan(B_j)}$ ($1\le j\le r$), add
\[
  A_{\out_\rho(\mathbf t_1,\ldots,\mathbf t_r)}
  \to
  \rho(B_{1,\mathbf t_1},\ldots,B_{r,\mathbf t_r}).
\]
Here \(B_{j,\mathbf t_j}\) denotes the typed copy of \(B_j\) indexed by
\(\mathbf t_j\).

Finally, the \emph{trimmed refinement} \(\widetilde G_0\) is obtained by
keeping only the typed nonterminals and typed rules that actually occur in some
successful derivation from \(\widetilde S\); the subscript $0$ denotes this
productive/reachable core.  From this point on, a ``typed nonterminal'' or
``typed rule'' means one of \(\widetilde G_0\) unless the full refinement
\(G^h\) is mentioned explicitly.
\end{definition}

\begin{proposition}[Yield-type invariants]
\label{prop:output-invariants}
For \(G^h\) and \(\widetilde G_0\):
\begin{enumerate}[label=(\roman*),leftmargin=*]
\item if \(A_{\mathbf t}\Rightarrow^*\tuple u\), then
\(h^{(\fan(A))}(\tuple u)=\mathbf t\), and erasing the type indices gives a
valid \(G\)-derivation of \(\tuple u\);
\item every terminal derivation
\(A\Rightarrow_G^*\tuple u\)
lifts uniquely to a derivation
\(A_{h^{(\fan(A))}(\tuple u)}\Rightarrow_{G^h}^*\tuple u.\)
In particular, every successful derivation \(S\Rightarrow_G^* w\) lifts
uniquely through the start rule
\(
  \widetilde S\to S_{(h(w))}
\);
\item \(L(\widetilde G_0)=L(G^h)=L(G)\).
\end{enumerate}
\end{proposition}

\begin{proof}
At a rule application, homomorphic evaluation of the rule template gives
\[
  h^{(\fan(A))}
  \bigl(\rho(\tuple u_1,\ldots,\tuple u_r)\bigr)
  =
  \out_\rho
  \bigl(h^{(\fan(B_1))}(\tuple u_1),\ldots,
        h^{(\fan(B_r))}(\tuple u_r)\bigr).
\]
The rank-zero case is immediate.  Induction on derivation height proves~(i),
and the same identity gives the unique terminal-derivation lifting in~(ii).  Erasing the
type indices gives one language inclusion.  Conversely, every successful
terminal derivation of \(G\) has the unique typed lifting described in~(ii),
including the uniquely determined start link
\(\widetilde S\to S_{(h(w))}\).  Finally, every successful typed derivation
consists entirely of rules retained by the trimming, so trimming preserves the
generated language.
\end{proof}

Yield typing changes only nonterminal labels, so every typed rule has the
same rule template as its untyped source.  Hence $\lambda$-freeness,
nondeletion, nonpermutation, and nonmerging are all inherited by the trimmed
refinement.

For a typed nonterminal $X=A_{\mathbf t}$, put
$L_X:=\{\tuple u:X\Rightarrow_{\widetilde G_0}^*\tuple u\}$.  Trimming makes
$L_X$ nonempty; if the target grammar is $\lambda$-free, every component
of every $\tuple u\in L_X$ is nonempty.

\begin{lemma}[Canonical contexts for typed nonterminals]
\label{lem:identity-canonical-contexts}
Assume the target grammar is good.  Every typed nonterminal $X$ of arity $d$ has a sentence context
$\chi(X)\in\mathcal E_d$ such that
$\chi(X)[\tuple u]\in L(G)$ for every $\tuple u\in L_X$.  If
$\widetilde S\to X$ belongs to $\widetilde G_0$, one may take $\chi(X)=\blank_1$.
\end{lemma}

\begin{proof}
Choose a successful typed derivation containing an occurrence of $X$.  At the
root, the one-hole context $\blank_1$ is valid.  Follow the unique path from
the root occurrence down to the chosen occurrence of $X$.  At each rule on
this path, fix the productive sibling subderivations used in the successful
derivation and apply Lemma~\ref{lem:filling-identity}.  Goodness of the rule
and nonemptiness of the sibling components ensure that the induced child
context again has nonempty internal separators.  Iterating down the path gives
$\chi(X)\in\mathcal E_d$.  The surrounding typed rule applications depend only on the typed symbol
$X$, so replacing its subtree by any other $X$-derivation remains a valid
derivation in $\widetilde G_0$.  If $\widetilde S\to X$ is a rule of
$\widetilde G_0$, the path has length zero and $\chi(X)=\blank_1$.
\end{proof}

When $L=L(G)\in\Cmcf{f}{h}$, the refinement is also distributionally homogeneous for the sentence-context
family used here.  If $X=A_{\mathbf t}$ has arity $d$ and
$\tuple u,\tuple v\in L_X$, Proposition~\ref{prop:output-invariants} gives
$h^{(d)}(\tuple u)=h^{(d)}(\tuple v)=\mathbf t$, while
Lemma~\ref{lem:identity-canonical-contexts} supplies one sentence context
accepting both.  Hence $(f,h)$-tuple substitutability gives
$\D_L^{(d)}(\tuple u)=\D_L^{(d)}(\tuple v)$.  Thus each typed nonterminal lies
inside one distributional class for $\mathcal E_d$.  This parallels, but is weaker than, the full-multicontext homogeneity used
for congruential MCFGs~\cite{ClarkYoshinaka2016}: our internal separators are
required to be nonempty.  For genuine tuple arity $d\ge2$, the present
sentence-context family still requires nonempty internal separators, so we do
not assert an inclusion into the full multicontext congruential MCFG class;
whether the empty-separator cases can be recovered uniformly remains open.

\begin{definition}[Canonical tuples and contexts]
\label{def:target-canonical}
For every non-start typed nonterminal $X$, fix a \emph{canonical tuple}
$\omega(X)\in L_X$ and a \emph{canonical context} $\chi(X)$ supplied by
Lemma~\ref{lem:identity-canonical-contexts}.  Whenever
$\widetilde S\to X$ belongs to $\widetilde G_0$, require
$\chi(X)=\blank_1$; otherwise choose any canonical context from the lemma.  For
a $\lambda$-free target grammar every component of $\omega(X)$ is nonempty,
so $\|\omega(X)\|_1\ge\fan(X)\ge1$.  Qualitatively a canonical tuple may be the first
tuple in an effective enumeration of finite terminal derivations; the
quantitative sections choose short canonical tuples.  These canonical objects are target witnesses, not inputs to the reconstruction
operator.
\end{definition}

\begin{definition}[Finite target witness set]
\label{def:witness-set}
Let $\Wit(\widetilde G_0)$ consist of the following positive words:
\[
\begin{aligned}
&\chi(X)[\omega(X)]
&&\text{for every non-start typed nonterminal }X,\\
&\chi(X)[\rho(\omega(Y_1),\ldots,\omega(Y_r))]
&&\text{for every non-start positive-rank typed rule }
  X\to\rho(Y_1,\ldots,Y_r),\\
&\chi(X)[\tuple c]
&&\text{for every non-start rank-zero typed rule }X\to\tuple c.
\end{aligned}
\]
\end{definition}

Because $\widetilde G_0$ is finite and all of its symbols and rules occur in
successful derivations, $\Wit(\widetilde G_0)$ is a finite subset of $L(G)$;
when $L(G)\ne\emptyset$ it is nonempty.

\begin{lemma}[Target rule words occur in the enumeration]
\label{lem:exposed-rule-enumerated}
Let $R:X\to\rho(Y_1,\ldots,Y_r)$ be a non-start positive-rank typed rule.
Put $\tuple z_R:=\rho(\omega(Y_1),\ldots,\omega(Y_r))$.  If
$\Wit(\widetilde G_0)\subseteq K$, the witness enumeration contains
\[
  [\tuple z_R]\to\rho([\omega(Y_1)],\ldots,[\omega(Y_r)]).
\]
\end{lemma}

\begin{proof}
The word $\chi(X)[\tuple z_R]\in K$ supplies the parent occurrence.
The segmentation of this sample occurrence induced by the variables of
$\rho$ places every canonical tuple of each child as a tuple occurrence in the same sample word.  $\lambda$-freeness makes all child components nonempty, nonpermutingness
puts each child's components in index order, and nonmerging guarantees a
nonempty separator between consecutive components of each child in the
sample word.  Thus every child canonical tuple is an observed tuple occurrence and the
segmentation is a composition witness; exhaustive enumeration emits the
displayed rule.
\end{proof}


\begin{proposition}[Cardinality of the finite target witness set]
\label{prop:witness-cardinality}
Let $G=(V,\Sigma,P,S,\fan)$ be the normalized target presentation from which
$\widetilde G_0$ is constructed, with fan-out at most
$f$ and rule rank at most $r_{\max}$.  For a fixed finite typing
$h:\Sigma^*\to M$, the target witness set constructed
above may be chosen so that
\[
  |\Wit(\widetilde G_0)|
  \le |V||M|^f+|P||M|^{f r_{\max}}.
\]
In particular, for fixed $f$, $r_{\max}$, and $h$, the number of selected positive
words is linear in $|V|+|P|$.
\end{proposition}

\begin{proof}
The yield-type refinement has at most $|V||M|^f$ non-start typed nonterminals.
The source rules give at most $|P||M|^{f r_{\max}}$ non-start typed rule
instances.  The construction selects at most one word for each such typed
nonterminal and one for each such typed rule instance; fresh start links are
handled separately and require no additional witness word.
\end{proof}

Proposition~\ref{prop:witness-cardinality} controls only the number of selected
words, not their total positive length $\|\Wit(\widetilde G_0)\|_+$.
Polynomial total-length bounds require the additional rank-one or
boundary-determined arguments below.


\subsection{Completeness and Exact Reconstruction}

\label{sec:completeness}

\begin{proposition}[Completeness from a target presentation]
\label{prop:completeness}
This is the finite-typed counterpart of Yoshinaka's rule-witness
completeness argument~\cite[Lemma~5]{Yoshinaka2011}.  Let $G$ be a finite good MCFG of
fan-out at most $f$, let
$L_G:=L(G)$, and construct its trimmed yield-type refinement
$\widetilde G_0$ and witness set above.  Let $r_G$ be the maximum rank of a
non-start typed rule, with $r_G:=0$ if there is none.  For every
finite sample $K$ with $\Wit(\widetilde G_0)\subseteq K$ and every
$r_{\max}\ge\max\{1,r_G\}$,
\[
  L_G\subseteq L\bigl(\mathcal B_{f,h}^{\le r_{\max}}(K)\bigr),
\]
and also $L_G\subseteq L\bigl(\mathcal B_{f,h}(K)\bigr)$.
\end{proposition}

\begin{proof}
For each non-start typed nonterminal $X$, put $\iota(X):=[\omega(X)]$.  If
$\widetilde S\to X$ belongs to $\widetilde G_0$, then $\chi(X)=\blank_1$, so the reconstruction contains
the corresponding start rule.

For a non-start rank-zero typed rule $R:X\to\tuple c$, the two corresponding witness words share $\chi(X)$, and Proposition~\ref{prop:output-invariants} gives the same
componentwise $h$-type to $\omega(X)$ and $\tuple c$.  Hence the reconstruction contains
$\iota(X)\to[\tuple c]$, followed by the constant rule.

For a non-start positive-rank typed rule
$R:X\to\rho(Y_1,\ldots,Y_r)$, put
$\tuple z_R:=\rho(\omega(Y_1),\ldots,\omega(Y_r))$.  Yield typing gives
$h^{(\fan(X))}(\tuple z_R)=h^{(\fan(X))}(\omega(X))$, and the two corresponding witness words share $\chi(X)$, so Rule~(Link)
$\iota(X)\to[\tuple z_R]$ is present.  Lemma~\ref{lem:exposed-rule-enumerated}
provides
\[
 [\tuple z_R]\to\rho(\iota(Y_1),\ldots,\iota(Y_r)).
\]
Induction on a $\widetilde G_0$-derivation therefore simulates every non-start
typed rule whenever the cap is at least $r_G$; the fresh start link was handled
above.  Hence every word of $L_G$ is generated.
For the uncapped reconstruction, each required rule has a witness word in $K$;
Lemma~\ref{lem:general-witness-enumeration} places its rank below $\|K\|_+$.  This
simulation uses neither tuple substitutability nor $K\subseteq L$; those enter
only through soundness.
\end{proof}

\begin{theorem}[Exact reconstruction for a fixed typing]
\label{thm:exact-reconstruction}
Fix $f\ge1$ and a finite homomorphism $h:\Sigma^*\to M$.  For every
$L\in\Cmcf{f}{h}$ there is a finite characteristic sample
$K_L^{\star}\subseteq L$ such that every finite
$K_L^{\star}\subseteq K\subseteq L$ satisfies
\[
  L(\mathcal B_{f,h}(K))=L.
\]
\end{theorem}

\begin{proof}
Put $L^+:=L\setminus\{\eps\}$.  If $L^+=\emptyset$, take $K_L^\star=L$.
Otherwise choose a finite $f$-MCFG presentation of $L$ and apply
Proposition~\ref{prop:good-normalization} to obtain a good presentation of
$L^+$.  Let $\widetilde G_0$ be its trimmed yield-type refinement and set
\[
 K_L^\star:=\Wit(\widetilde G_0)\cup(\{\eps\}\cap L).
\]
The observation following Definition~\ref{def:witness-set} makes this a finite
subset of $L$.  For $K_L^\star\subseteq K\subseteq L$,
Proposition~\ref{prop:completeness} gives every nonempty target word in the
hypothesis, while Proposition~\ref{prop:reconstruction-soundness} gives the
reverse inclusion.  If $\eps\in L$, the direct start rule generates it; if
$\eps\notin L$, soundness forbids it.  The cases $L=\emptyset$ and
$L=\{\eps\}$ are immediate from the reconstruction definition.
\end{proof}

The set-driven map that sends a finite text prefix to
$\mathcal B_{f,h}$ of its positive content is eventually semantically correct:
once the content contains $K_L^\star$, every later reconstruction generates
$L$ by Theorem~\ref{thm:exact-reconstruction}.  Its grammar presentation may
still change as new positive examples enlarge the observed tuple inventory,
which is why TxtEx stabilization is handled by the conservative wrapper below.

\begin{definition}[Conservative sequential learner]
\label{def:conservative-learner}
We use the same conservative update policy as Yoshinaka's
Algorithm~1~\cite[Section~4.1]{Yoshinaka2011}.  Let $H_0$ be the fixed
empty-language hypothesis from the learning model.  The
learner $\mathcal A_{f,h}$ starts from $H_0$.  On a pause symbol it keeps its
current grammar.  On a positive datum $w$, it keeps the current grammar if that
grammar already generates $w$; otherwise, with $K$ the positive content seen
so far, it replaces the hypothesis by $\mathcal B_{f,h}(K)$.
\end{definition}

\begin{corollary}[Gold/TxtEx identification]
\label{cor:gold-explanatory}
For every finite-monoid homomorphism $h:\Sigma^*\to M$ and every $f\ge1$,
the conservative learner $\mathcal A_{f,h}$ TxtEx-identifies
$\Cmcf{f}{h}$ from positive data.
\end{corollary}

\begin{proof}
MCFG membership is decidable~\cite{SekiEtAl1991}, so the learner is
computable.  Sample consistency gives the invariant that every current
hypothesis contains every positive example seen so far: a rebuild contains the
whole current content, and a stage without a rebuild leaves a hypothesis that
already contains the new datum.  Proposition~\ref{prop:reconstruction-soundness}
makes every rebuilt hypothesis sound.

Fix a stage after the characteristic sample from
Theorem~\ref{thm:exact-reconstruction} has appeared.  If the current sound
hypothesis is already exact, every later positive datum is generated and no
further rebuild occurs.  Otherwise choose $w\in L$ missing from the current
hypothesis.  The text eventually presents $w$, which forces a rebuild from
positive content still containing the characteristic sample; that rebuild is
exact by Theorem~\ref{thm:exact-reconstruction}.  From then on no further
change is possible.  The learner is generally presentation-order-sensitive and
need not be set-driven.
\end{proof}

\begin{proposition}[Unknown typing with a bounded codomain]
\label{prop:bounded-unknown-typing}
Fix $\Sigma$, $f$, and $m_{\max}\ge1$, and let
$\Cboundedtyping{f}{m_{\max}}$ be the union of $\Cmcf{f}{h}$ over all
homomorphisms $h:\Sigma^*\to M$ with $|M|\le m_{\max}$.  There is one finite
homomorphism $U_{\Sigma,m_{\max}}$ such that
\[
 \Cboundedtyping{f}{m_{\max}}
 \subseteq \Cmcf{f}{U_{\Sigma,m_{\max}}},
\]
so the bounded union is TxtEx-identifiable.  In contrast,
$\Calltyping{f}=\bigcup_h\Cmcf{f}{h}$ is not TxtEx-identifiable from positive
data.
\end{proposition}

\begin{proof}
There are only finitely many monoid structures of size at most $m_{\max}$ and
letter maps from $\Sigma$ into them.  Take the product of all resulting
homomorphisms and restrict it to its image.  Every admissible $h$ factors
through this product, so typing refinement gives the displayed inclusion and
Corollary~\ref{cor:gold-explanatory} gives TxtEx identification.  The product
can be extremely large; no polynomial dependence on $m_{\max}$ is claimed.
Without a codomain bound the union contains every regular language, hence all
finite languages and an infinite language, so Gold's obstruction for a class containing all finite languages together with an
infinite language applies~\cite{Gold1967}.
\end{proof}

\subsection{Polynomial-Time Hypothesis Construction at Bounded Rule Rank}
\label{sec:poly-time}

The uncapped reconstruction allows arbitrary finite rule rank and serves the
qualitative identification result; polynomial update bounds are stated only
for the capped regime.  When $f$, $r_{\max}$, and $h$ are treated as fixed
class parameters, the same witness construction is polynomial in the observed
sample.  This is the standard bounded-dimension/bounded-rank regime used for
polynomial MCFG learning and recognition; Yoshinaka likewise fixes dimension
and rule rank in his polynomial learner and invokes the corresponding uniform
membership bound \cite[Proposition~4 and Section~4.3]{Yoshinaka2011}.

\paragraph{Encoding convention.}
We use canonical finite encodings for all sample-derived objects.  An occurrence is
recorded by its sample-word index and cut positions;
tuple contexts and templates are recorded by arity together with delimited
terminal factors and ordered hole or variable labels.

\begin{theorem}[Polynomial construction under a rule-rank cap]
\label{thm:poly}
Fix a fan-out bound \(f\ge1\), a rule-rank cap \(r_{\max}\ge1\), and a
finite-monoid homomorphism \(h:\Sigma^*\to M\).  From any finite positive
sample \(K\), the capped reconstruction hypothesis
\(\widehat G_{\le r_{\max}}(K)\)
can be constructed, including its complete output encoding, in time
\((1+\|K\|_+)^{O(f r_{\max})}.\)
Its fan-out is at most \(f\), and every composition-witness rule has rank at
most \(r_{\max}\).
\end{theorem}

Rule~(Link) is unary and Rule~(Const) has rank zero.  Here $h$ is fixed: hidden constants may depend on its finite
multiplication table and letter images, while the polynomial degree may depend
on $f$ and $r_{\max}$.  The complexity statement treats $h$ as a fixed class parameter; a variable
input encoding of $h$ lies outside its scope.

\begin{proof}
Put \(n_K:=\max(1,\|K\|_+)\).  Tuple occurrences of arity at most \(f\) are
enumerable in \(n_K^{O(f)}\) time by the same cut-position count used in
Lemma~\ref{lem:general-witness-enumeration}.  All observed constant rules are
read from these tuples.  Since $h$ is given explicitly, all substring $h$-values in
the sample can be precomputed in polynomial time (for example, by extending
from every cut position), so type evaluation introduces no hidden
super-polynomial cost.  Rule~(Link) is obtained by grouping
occurrences by arity and concrete-context encoding and comparing componentwise
\(h\)-types; even exhaustive pairwise comparison remains within
\(n_K^{O(f)}\).

By Lemma~\ref{lem:general-witness-enumeration}, all composition
witnesses of rank at most \(r_{\max}\), together with their canonical template
encodings, are enumerable in \(n_K^{O(f r_{\max})}\) time and space.  Canonical sorting
removes duplicate nonterminals and rules within the same bound.  Thus the
number of output symbols, rules, and their total encoding length are all
\(n_K^{O(f r_{\max})}\).
\end{proof}

\begin{corollary}[Polynomial-time conservative learning at bounded rank]
\label{cor:capped-gold-poly}
Fix $f\ge1$, $r_{\max}\ge1$, and a homomorphism
$h:\Sigma^*\to M$ into a finite monoid.  There is a conservative learner
$\mathcal A_{f,h}^{\le r_{\max}}$ that TxtEx-identifies
\[
  \Cmcf{f}{h}\cap \mathsf{MCFL}(f,r_{\max})
\]
and processes each datum in time polynomial in the accumulated positive-sample
size, which includes the length of the newly received datum.  Whenever it changes hypothesis, it rebuilds
$\mathcal B_{f,h}^{\le r_{\max}}(K)$ from the current content $K$.
\end{corollary}

\begin{proof}
Use the policy of Definition~\ref{def:conservative-learner}, but replace the
uncapped reconstruction by $\mathcal B_{f,h}^{\le r_{\max}}$.  For a target in
$\mathsf{MCFL}(f,r_{\max})$, choose a rank-$\le r_{\max}$ presentation,
apply the rank-preserving good normalization, and form the characteristic
sample of Theorem~\ref{thm:exact-reconstruction} from that normalized
presentation.  Proposition~\ref{prop:completeness} then makes every rebuild
after this sample exact; soundness and sample consistency are unchanged.
Hence the stabilization proof of Corollary~\ref{cor:gold-explanatory} applies
verbatim.

By Theorem~\ref{thm:poly}, a rebuild has polynomial cost for fixed
$f,r_{\max},h$.  Every current hypothesis has fan-out at most $f$ and rule
rank at most $r_{\max}$: Rules~(Start)--(Const) have rank zero, Rule~(Comp) is capped,
and Rule~(Link) is unary.  For fixed fan-out and rule rank, MCFG membership is polynomial jointly in
the grammar encoding and input length
\cite[Proposition~4]{Yoshinaka2011}; see also
\cite{SekiEtAl1991,Kallmeyer2010} for the classical recognition and parsing
background.  Unary Rule~(Link) cycles cause no
difficulty: the chart algorithm saturates the finite unary closure over the
current nonterminal set.  Thus both the membership test that decides whether
to rebuild and the rebuild itself take polynomial time in the accumulated
positive-sample size.
\end{proof}

At fan-out one, every context-free language has a rule-rank-two CFG
presentation after the usual Chomsky-normal-form conversion, with the empty
word handled separately at the start symbol.  Hence, for every fixed finite
$h$, the capped learner $\mathcal A_{1,h}^{\le2}$ already covers the whole
class $\Cmcf{1}{h}$ and has polynomial time per processed datum.  No grammar-size-only characteristic-data bound is implied at rank two.

\section{Characteristic-Data Bounds}
\label{sec:char-data}

Under a rule-rank cap, Proposition~\ref{prop:witness-cardinality} already
gives polynomially many characteristic words; this section controls their
total length.

\begin{center}
\small
\begin{tabularx}{\textwidth}{>{\raggedright\arraybackslash}p{0.27\textwidth}
                              >{\raggedright\arraybackslash}p{0.35\textwidth}
                              X}
\hline
Condition & Characteristic-data control & Source \\
\hline
arbitrary $h$, bounded rank
  & polynomial in $\|G\|$ and refined thickness $\theta_{\widetilde G_0}$
  & Lemma~\ref{lem:typed-thickness-baseline} \\
rule rank one
  & polynomial in $\|G\|$ alone
  & Theorem~\ref{thm:rankone-polydata} \\
auxiliary: unique componentwise $h$-type for every rule child
  & explicit $O((|V|+|P|)(1+r_{\max})\max\{1,\theta_G\})$ scale
  & Corollary~\ref{cor:no-splitting-thickness} \\
$h\preceq h_\Sigma^{k,\ell}$, bounded rank
  & polynomial in $\|G\|$ and $\max\{1,\theta_G\}$
  & Theorem~\ref{thm:boundary-polydata} \\
\hline
\end{tabularx}
\end{center}

The first row is the unconditional baseline for general $h$.
Full yield-type refinement can make its thickness exponentially larger than
the original presentation parameters
(Proposition~\ref{prop:typed-thickness-exponential}), so that route alone
gives no original-thickness bound.  Rank one removes the thickness dependence,
whereas prefix--suffix-determined typing gives the main higher-rank transfer;
the unique-child-type row is only an auxiliary no-splitting criterion.
Polynomial update time in the observed sample is independent of these
data-length refinements
(Theorem~\ref{thm:poly} and Corollary~\ref{cor:capped-gold-poly}).

\subsection{Thickness and reusable canonical contexts}
\label{subsec:thickness-contexts}

\paragraph{Rule thickness and Yoshinaka's efficiency criterion.}
For the quantitative results, presentations are reduced and \emph{good}, as in
Yoshinaka's fixed-dimension, fixed-rank good form
\cite[Section~2.3, after Example~2]{Yoshinaka2011}.  For a rule $R$ of $G$ let
$\theta_G(R)$ be the minimum length of a successful sentence derivation using
$R$, put $\theta_G:=\max_{R\in P}\theta_G(R)$, and define
\[
 \tau_G:=\max_A\min\{\|\tuple w\|_1:\tuple w\in L_A(G)\}.
\]
Thus $\theta_G$ is rule thickness and $\tau_G$ nonterminal thickness.  In
Yoshinaka's notation these correspond respectively to his $\tau_G$ and $t_G$,
with $\tau_G\le\|G\|t_G$ in his notation
\cite[Section~4.3]{Yoshinaka2011}.

\begin{lemma}[Polynomial equivalence of the two thickness parameters]
\label{lem:thickness-equivalence}
For every reduced good presentation $G$,
\[
  \tau_G\le\theta_G
  \qquad\text{and}\qquad
  \theta_G\le 2\|G\|\max\{1,\tau_G\}.
\]
\end{lemma}

\begin{proof}
For $\tau_G\le\theta_G$, choose any rule with left-hand side $A$ and place it
in a successful derivation.  Nondeletion preserves the tuple produced at that
occurrence of $A$, so its total length is bounded by the final sentence.
Minimize for $A$ and maximize over $A$.

For the converse, fix a rule $R$ with left-hand side $A$ and follow a shortest
dependency path from the start symbol to $A$.  The path repeats no
nonterminal.  Keep the path child open, complete each sibling by a shortest
terminal tuple, then apply $R$ and complete its children likewise.  Across the
distinct rules on this construction, both explicit terminal material and
child-variable occurrences are bounded by $\|G\|$, giving a successful
sentence of length at most
$\|G\|+\|G\|\tau_G\le2\|G\|\max\{1,\tau_G\}$.
\end{proof}

Hence the two thickness parameters are polynomially equivalent on reduced good
presentations.  Yoshinaka's efficiency criterion asks for polynomial
hypothesis construction, characteristic-set cardinality, and total positive
data in target size and rule thickness
\cite[Definition~1]{Yoshinaka2011}; see also
\cite{EyraudHeinzYoshinaka2016}.  Here characteristic-data bounds concern the
capped set-driven reconstruction, while TxtEx stabilization is supplied by the
conservative learner.  Proposition~\ref{prop:structural-good-thickness}
transfers the size--thickness scale for reduced $\lambda$-free inputs; genuine
empty components may instead cause the exponential growth of
Proposition~\ref{prop:empty-component-blowup}.

\begin{lemma}[Reusable canonical-context bound]

\label{lem:reusable-context-bound}
Let $\mathcal G_{\mathrm t}$ be a trimmed good typed MCFG
with $n_{\mathrm{typ}}$ non-start typed nonterminals and maximum non-start rule rank at most
$r_0\ge1$.  Suppose that every non-start typed nonterminal $Y$ has a canonical tuple
$\omega(Y)\in L_Y(\mathcal G_{\mathrm t})$ with
$\|\omega(Y)\|_1\le b_{\mathrm{can}}$, and that every non-start typed rule writes at
most $b_{\mathrm{rule}}$ terminal symbols explicitly.  Then the
canonical contexts may be chosen so that
\[
  |\chi(X)|_{\mathrm{fix}}
  \le n_{\mathrm{typ}}\bigl(b_{\mathrm{rule}}+(r_0-1)b_{\mathrm{can}}\bigr)
\]
for every non-start typed nonterminal $X$.  With these contexts and canonical tuples,
every selected witness word has length at most
\[
  n_{\mathrm{typ}}\bigl(b_{\mathrm{rule}}+(r_0-1)b_{\mathrm{can}}\bigr)+b_{\mathrm{rule}}+r_0b_{\mathrm{can}}.
\]
\end{lemma}

\begin{proof}
If $\widetilde S\to X$ is present take $\chi(X)=\blank_1$; otherwise follow a
shortest typed-nonterminal dependency path to $X$.  It has at most
$n_{\mathrm{typ}}$ non-start steps.  Keeping the path child open and filling
each sibling canonically costs at most
$b_{\mathrm{rule}}+(r_0-1)b_{\mathrm{can}}$ terminals per step.  Goodness
preserves the selected components, their order, and nonempty separators,
giving the context bound.  Adding either one canonical tuple or one rule
template with at most $r_0$ child tuples gives the displayed witness bound.
\end{proof}

\subsection{Refined-Thickness Baseline}
\label{subsec:typed-thickness}

For a reduced good presentation $G$ of an $\eps$-free target, let
$\widetilde G_0$ be its trimmed full yield-type refinement and put
\[
 \bar\theta_{\mathrm{typ}}
 :=\max\{1,\theta_{\widetilde G_0}\},
\]
where $\theta_{\widetilde G_0}$ is rule thickness computed in the refined
grammar.  No comparison with the original untyped thickness is assumed.

\begin{lemma}[Refined-thickness baseline]
\label{lem:typed-thickness-baseline}
Fix $f\ge1$, $r_{\max}\ge1$, a finite alphabet $\Sigma$, and a finite
typing $h:\Sigma^*\to M$.  Let $G$ be a reduced good fan-out-$\le f$,
rank-$\le r_{\max}$ presentation of an $\eps$-free language in
$\Cmcf{f}{h}$, and let $\widetilde G_0$ be its trimmed yield-type refinement.
Then the target witness set has cardinality polynomial in $\|G\|$ and total
positive length polynomial in $\|G\|$ and
$\bar\theta_{\mathrm{typ}}:=\max\{1,\theta_{\widetilde G_0}\}$.
\end{lemma}

\begin{proof}
The trimmed refinement $\widetilde G_0$ remains reduced and good.  A
minimum-yield successful derivation using a typed rule with left-hand side
$X$ contains an $X$-tuple of total length at most
$\bar\theta_{\mathrm{typ}}$, so use such tuples canonically.  If
$\widetilde S\to X$ exists take $\chi(X)=\blank_1$; otherwise remove the
chosen $X$-subtree from the same derivation to obtain a canonical context with
fixed material at most $\bar\theta_{\mathrm{typ}}$.  Each typed rule writes at
most that many terminals and has at most $r_{\max}$ canonical children, so
every selected witness has length at most
$(r_{\max}+2)\bar\theta_{\mathrm{typ}}$.  Therefore
\[
 |\Wit(\widetilde G_0)|
 \le |V||M|^f+|P||M|^{fr_{\max}},
 \qquad
 \|\Wit(\widetilde G_0)\|_+
 \le
 (|V||M|^f+|P||M|^{fr_{\max}})
 (1+(r_{\max}+2)\bar\theta_{\mathrm{typ}}),
\]
which is polynomial for fixed $f,r_{\max},h$.
\end{proof}

\begin{corollary}[Auxiliary no-splitting criterion]
\label{cor:no-splitting-thickness}
Under the assumptions of Lemma~\ref{lem:typed-thickness-baseline}, suppose in
addition that for every nonterminal $B$ occurring as a child on the right-hand
side of a rule there is a tuple
$\mathbf t_B\in M^{\fan(B)}$ such that
\[
  h^{(\fan(B))}(\tuple u)=\mathbf t_B
  \qquad\text{for every }\tuple u\in L_B(G).
\]
Then every source rule of $G$ has exactly one typed instance in
$\widetilde G_0$ and
\[
  \theta_{\widetilde G_0}=\theta_G.
\]
Moreover,
\[
  |\Wit(\widetilde G_0)|\le |V|+2|P|,
\]
and the witness set may be chosen so that
\[
  \|\Wit(\widetilde G_0)\|_+
  \le
  (|V|+2|P|)
  \bigl(1+(r_{\max}+2)\max\{1,\theta_G\}\bigr).
\]
\end{corollary}

\begin{proof}
Each productive child has its unique type $\mathbf t_B$; hence every productive
source rule has one typed instance, whose parent type is forced by
$\out_\rho$.  Reducedness and unique lifting show that these instances survive
trimming, and erasing/reinstating type indices gives a yield-length-preserving
bijection between successful uses of a source rule and its typed instance.
Thus $\theta_{\widetilde G_0}=\theta_G$ (start links have no larger
thickness).  There is at most one retained typed copy of each ordinary child
nonterminal, at most $|P|$ typed copies of the original start symbol, and one
typed instance per source rule, giving at most $|V|+2|P|$ witnesses.
Lemma~\ref{lem:typed-thickness-baseline} then gives the stated length bound.
\end{proof}

The singleton-type hypothesis is effectively checkable: the finite sets
$h^{(\fan(A))}(L_A(G))$ are obtained by saturation using the template maps
$\out_\rho$, polynomially for fixed $f,r_{\max},h$.  Outside this regime,
refined thickness need not be controlled by the source presentation; rank one
and prefix--suffix-determined typing are the two positive cases below.

\begin{proposition}[Exponential blow-up of full yield-type refinement]
\label{prop:typed-thickness-exponential}
There are a fixed two-element typing $h_{\mathtt c}$, reduced good fan-out-one rank-two
grammars $G_n$ of size $O(n)$, and a constant bound
$\theta_{G_n}\le2$, such that $L(G_n)=\{\mathtt a,\mathtt c\}^+$ and the
trimmed full yield-type refinement $\widetilde G_{n,0}$ satisfies
\[
  \theta_{\widetilde G_{n,0}}\ge 2^n.
\]
\end{proposition}

\begin{proof}
Let $M_{\mathtt c}=\{1,\zeta\}$ with $\zeta^2=\zeta$,
$h_{\mathtt c}(\mathtt a)=1$, and
$h_{\mathtt c}(\mathtt c)=\zeta$.  For $n\ge1$, view
\[
\begin{aligned}
 S&\to SS\mid\mathtt a\mid\mathtt c\mid A_0\mid\cdots\mid A_n,\\
 A_0&\to\mathtt a,\qquad
 A_i\to A_{i-1}A_{i-1}\mid\mathtt c\quad(1\le i\le n)
\end{aligned}
\]
as a fan-out-one MCFG.  It is reduced and good, has size $O(n)$ and rule rank
two, generates $\{\mathtt a,\mathtt c\}^+$, and every rule has a successful
use of length at most two, so $\theta_{G_n}\le2$.

In the full refinement, the retained copy $(A_n)_1$ can derive only words
without $\mathtt c$; recursively the unique such yield is
$\mathtt a^{2^n}$.  Hence the typed rule
$S_{(1)}\to(A_n)_1$ has successful thickness at least $2^n$.
\end{proof}

Thus full yield-type refinement admits no bound polynomial in
$\|G\|\max\{1,\theta_G\}$ for arbitrary fixed typing, already at fan-out one
and rank two.  This is a limitation of the full-refinement route, not of the
language $\{\mathtt a,\mathtt c\}^+$ itself.

\subsection{Rule Rank One: Grammar-Size Polynomial Data}
\label{subsec:rankone-polydata}

At rule rank one every terminal derivation in the typed refinement is a unary
chain ending in a rank-zero rule, so its length can grow only by the terminals
written along that chain.  Let $n_{\mathrm{typ}}$ be the number of non-start
typed nonterminals of $\widetilde G_0$ and let $b_G\ge1$ be the maximum number
of terminals explicitly written by a rule.  Then
\[
 n_{\mathrm{typ}}\le |V||M|^f,\qquad b_G\le\max\{1,\|G\|\}.
\]

\begin{lemma}[Short typed canonical tuples at rule rank one]
\label{lem:rankone-short-canonical}
For every non-start typed nonterminal $X$ of $\widetilde G_0$, the
canonical tuple may be chosen so that
\[
 \|\omega(X)\|_1\le n_{\mathrm{typ}}b_G.
\]
\end{lemma}

\begin{proof}
Choose a shortest terminal unary chain from $X$.  It cannot repeat a typed
nonterminal, since splicing in the lower suffix would shorten the chain.
Thus it has at most $n_{\mathrm{typ}}$ rule applications, each adding at most
$b_G$ terminals because the rules are linear and nondeleting.
\end{proof}

Lemma~\ref{lem:reusable-context-bound} with
$b_{\mathrm{can}}=n_{\mathrm{typ}}b_G$, $b_{\mathrm{rule}}=b_G$, and $r_0=1$
then bounds every selected characteristic word by
$(2n_{\mathrm{typ}}+1)b_G$.

\begin{theorem}[Rank-one polynomial construction and characteristic data]
\label{thm:rankone-polydata}
Fix $f\ge1$, a finite alphabet $\Sigma$, and a finite typing
$h:\Sigma^*\to M$.  Let $G_{\mathrm{in}}$ be any finite standard fan-out-$\le f$,
rule-rank-$\le1$ presentation of a target $L\in\Cmcf{f}{h}$.  Then
$\mathcal B_{f,h}^{\le1}$ is computable in time polynomial in $\|K\|_+$ for a finite positive
sample $K$ and admits characteristic data of total positive length
$\operatorname{poly}_{f,h}(\|G_{\mathrm{in}}\|)$.  If $\eps\in L$, the
characteristic sample additionally contains the single direct-start datum
$\eps$; its length is zero and hence does not affect this bound.  The capped conservative learner
$\mathcal A_{f,h}^{\le1}$ TxtEx-identifies
$\Cmcf{f}{h}\cap\mathsf{MCFL}(f,1)$, with polynomial time per processed datum.
\end{theorem}

\begin{proof}
A rank-zero rule has one retained typed instance and a unary rule at most
$|M|^f$, so the refined rule set has size $O(|P||M|^f)$.  By the preceding
lemma and Proposition~\ref{prop:witness-cardinality},
\[
 \|\Wit(\widetilde G_0)\|_+
 \le (n_{\mathrm{typ}}+|P||M|^f)
      (1+(2n_{\mathrm{typ}}+1)b_G),
 \qquad n_{\mathrm{typ}}\le |V||M|^f.
\]
This is polynomial for fixed $f,h$.  Proposition~\ref{prop:good-normalization}
preserves rank one with fixed-parameter linear overhead, while
Theorem~\ref{thm:poly} and Corollary~\ref{cor:capped-gold-poly} give the stated
reconstruction and conservative TxtEx bounds.
\end{proof}

At $f=1$ this is the linear-CFG unary-chain argument.  Rank two already
precludes a grammar-size-only data bound: the size-$O(n)$ grammar
$A_0\to\mathtt c$, $A_i\to A_{i-1}A_{i-1}$, $S\to A_n$ generates the singleton
$\{\mathtt c^{2^n}\}$, so any nonempty characteristic sample has length at
least $2^n$.  This motivates the thickness parameter used below.

\subsection{Prefix--Suffix-Determined Typings: Thickness-Sensitive Data}
\label{subsec:boundary-polydata}

Lemma~\ref{lem:typed-thickness-baseline} gives an arbitrary-$h$
characteristic-data bound in terms of the thickness of the refined typed
presentation.  Proposition~\ref{prop:typed-thickness-exponential} shows that
this refined thickness need not be polynomially controlled by the original
good presentation.  A stronger transfer is nevertheless possible whenever the
finite typing is determined by a fixed prefix--suffix window.  We formulate
this boundary-determined regime in the same thickness-sensitive representation
model used by Yoshinaka~\cite[Section~4.3]{Yoshinaka2011}.

Let $G=(V,\Sigma,P,S,\fan)$ be a reduced good MCFG and put
\[
  \bar\theta_G:=\max\{1,\theta_G\}.
\]
Here $\theta_G$ is the rule-thickness parameter defined above.  The argument is stated for the
$\lambda$-free good core.  If the target also contains $\eps$, apply
the argument to $L\setminus\{\eps\}$ and add the single datum $\eps$ at
the start symbol, exactly as in Theorem~\ref{thm:exact-reconstruction}; this
adds only one length-zero example and does not affect either polynomial bound.

\begin{lemma}[Rule thickness bounds local terminal material]
\label{lem:thickness-local-bounds}
Let $G$ be reduced and good.  Then:
\begin{enumerate}[label=(\roman*),leftmargin=*]
\item for every rule $R$, the total number of terminal symbols written
explicitly in the template of $R$ is at most $\theta_G$;
\item for every nonterminal $A$, there exists a terminal tuple
$\tuple u\in L_A(G)$ with $\|\tuple u\|_1\le\theta_G$.
\end{enumerate}
\end{lemma}

\begin{proof}
Reducedness places every rule in some successful derivation.  Because a good
grammar is nondeleting, every terminal written by that rule and every component
of the tuple produced at its left-hand side survives exactly once to the final
sentence.  A minimum-yield successful derivation using the rule therefore bounds both
its explicit terminal material and the total length of that occurrence's
left-hand-side tuple.  Maximizing over rules gives (i) and (ii).
\end{proof}

Fix $k,\ell\ge0$ and let
$h:\Sigma^*\to M$ be any finite typing with
$h\preceq h_\Sigma^{k,\ell}$.  Thus $h$ factors through the surjective prefix--suffix typing
$h_\Sigma^{k,\ell}$.  Let $G$ have fan-out at most $f$ and
rule rank at most a fixed $r_{\max}\ge1$, and let $\widetilde G_0$ be its
trimmed full yield-type refinement under $h$.  Let $n_{\mathrm{typ}}$ be the
number of retained non-start $h$-typed nonterminals.

\begin{lemma}[Short prefix--suffix-typed canonical tuples]
\label{lem:boundary-short-canonical}
For every non-start typed nonterminal $X$ of $\widetilde G_0$, the
canonical tuple may be chosen so that
\[
  \|\omega(X)\|_1
  \le f(k+\ell+1)n_{\mathrm{typ}}(r_{\max}+1)\bar\theta_G.
\]
\end{lemma}

\begin{proof}
Choose, among all terminal derivations from $X$ in the trimmed typed grammar
$\widetilde G_0$, one with the minimum number of grammar-rule nodes, and then mark the root terminals needed to
determine its componentwise $h_\Sigma^{k,\ell}$-type: all terminals of a
component shorter than $m_0=\max\{1,k+\ell\}$ from
Definition~\ref{def:prefix-suffix-typing}, and otherwise its first $k$ and
last $\ell$ terminals (one arbitrary terminal when $k=\ell=0$).  There are at most $f(k+\ell+1)$ marked terminals.

No typed nonterminal can occur twice on a root-to-leaf grammar path of this minimal
derivation.  If an upper and a proper lower occurrence had the same typed
label, the lower subtree would itself be a valid typed terminal derivation from
the same typed nonterminal, and replacing the upper subtree by it would strictly reduce
the number of grammar-rule nodes.  Every typed nonterminal occurring in this derivation is retained in $\widetilde G_0$; hence every root-to-marked-terminal path
contains at most $n_{\mathrm{typ}}$ grammar nodes.  Their union therefore contains at most
$f(k+\ell+1)n_{\mathrm{typ}}$ grammar nodes.

Keep the union of these marked root-to-leaf paths as the retained subderivation.  By linearity and nondeletion, each component
of every omitted child contributes a contiguous interval to exactly one root
component; it is never split between root components.  Replace every omitted
child subtree by a shortest terminal tuple of its underlying untyped
nonterminal.  Lemma~\ref{lem:thickness-local-bounds}(ii) bounds each such tuple
by $\bar\theta_G$, and part~(i) bounds the explicit terminal material of each
retained rule by $\bar\theta_G$.  Because the grammar is good, every
replacement component is nonempty.  An omitted child cannot contribute to a
short root component: every terminal of such a component was marked, so any
nonempty contribution from that child would contain a marked terminal and the
corresponding root-to-terminal path would place the child on the retained
skeleton.  In a long root component no replaced block crosses a protected
prefix or suffix position.
Hence every replacement lies wholly in the unmarked interior of that component:
the marked first $k$ and last $\ell$ terminals remain literal prefixes and
suffixes, and the component stays on the long-word side of the threshold.
Thus the root prefix--suffix type is unchanged (and when $k=\ell=0$ the single
marked terminal preserves nonemptiness).  Applying the factor map supplied by
$h\preceq h_\Sigma^{k,\ell}$ shows that its componentwise $h$-type is
unchanged as well.
Each retained node contributes at most $\bar\theta_G$ explicit terminals and
has at most $r_{\max}$ omitted child subtrees.  The resulting root tuple has
length at most
\[
  f(k+\ell+1)n_{\mathrm{typ}}(r_{\max}+1)\bar\theta_G.
\]
It therefore has the $h$-type indexing $X$.  The replacements may change
internal typed labels of the original typed derivation; this is harmless.  We
retain only the underlying $G$-derivation and the preserved root type, and then
Proposition~\ref{prop:output-invariants}(ii) uniquely lifts that derivation back
to $G^h$ with root $X$.  Since $X$ belongs to $\widetilde G_0$, the lifted
derivation can be inserted into a successful derivation through $X$ and hence
survives trimming.  Its root tuple may therefore be chosen as $\omega(X)$.
\end{proof}

Applying Lemma~\ref{lem:reusable-context-bound} with
\[
 b_{\mathrm{can}}
 =f(k+\ell+1)n_{\mathrm{typ}}(r_{\max}+1)\bar\theta_G,
 \qquad
 b_{\mathrm{rule}}=\bar\theta_G,
 \qquad
 r_0=r_{\max},
\]
shows directly that every selected witness word has length
$O_{f,r_{\max},k,\ell,h,\Sigma}(\|G\|^2\bar\theta_G)$, since
$n_{\mathrm{typ}}\le |V||M|^f$.

\begin{theorem}[Polynomial data for prefix--suffix-determined typings]
\label{thm:boundary-polydata}
Fix $f\ge1$, $r_{\max}\ge1$, $k$, $\ell$, a finite alphabet $\Sigma$, and a
finite typing $h:\Sigma^*\to M$ satisfying
$h\preceq h_\Sigma^{k,\ell}$.  Let $G$ be a reduced good fan-out-$\le f$,
rank-$\le r_{\max}$ presentation of a language in $\Cmcf{f}{h}$.  Then
$\mathcal B_{f,h}^{\le r_{\max}}$ is polynomial-time in the positive sample.
A characteristic sample can be chosen with
\[
 |\Wit(\widetilde G_0)|=O_{f,r_{\max},h}(\|G\|),
 \qquad
 \|\Wit(\widetilde G_0)\|_+
 =O_{f,r_{\max},k,\ell,h,\Sigma}
   (\|G\|^3\bar\theta_G),
\]
where $\bar\theta_G=\max\{1,\theta_G\}$.  For a reduced $\lambda$-free input
presentation $G_{\mathrm{in}}$ (deleting rules allowed), the structural
good-form conversion of Proposition~\ref{prop:structural-good-thickness}
yields total data polynomial in
$\|G_{\mathrm{in}}\|\max\{1,\tau_{G_{\mathrm{in}}}\}$.
\end{theorem}

\begin{proof}
The $h$-refinement has at most $|V||M|^f$ typed nonterminals and
$|P||M|^{fr_{\max}}$ typed rule instances.
Lemma~\ref{lem:boundary-short-canonical} applies because preserving the finer
$h_\Sigma^{k,\ell}$ boundary type preserves the coarser $h$-type.  Together
with Lemma~\ref{lem:reusable-context-bound}, it bounds every selected witness
word by
$O_{f,r_{\max},k,\ell,h,\Sigma}(\|G\|^2\bar\theta_G)$.
Proposition~\ref{prop:witness-cardinality} gives
\[
 |\Wit(\widetilde G_0)|
 \le n_{\mathrm{typ}}+|P||M|^{f r_{\max}},
\]
with $n_{\mathrm{typ}}\le |V||M|^f$.  Multiplying the polynomial number of
witnesses by the preceding word-length bound, and adding at most one datum for
$\eps$, proves the asserted total-positive-length bound.  The same successful
witnesses also imply a polynomial bound on the refined typed rule thickness,
but that intermediate parameter is not needed for the estimate.  Theorem~\ref{thm:poly} gives the hypothesis-construction claim.
\end{proof}

The input-presentation clause follows from
Proposition~\ref{prop:structural-good-thickness}; nonterminal thickness is
used because deleting source rules may erase material before good-form
conversion.  Corollary~\ref{cor:capped-gold-poly} then gives polynomial time
per processed datum for the corresponding capped conservative learner.

The bound is relative to a reduced good bounded-rank presentation.
Theorem~\ref{thm:intrinsic-ranktwo-witness} below shows that this higher-rank
regime is needed at the language level: it gives a non-context-free
$(1,1)$-typed target of minimum fan-out two and minimum rule rank two even
when competing presentations may use arbitrary finite fan-out.

\section{Comparison with Multidimensional Substitutability}
\label{sec:comparison}

Yoshinaka's multidimensional substitutability is the direct positive-data
baseline~\cite{Yoshinaka2009,Yoshinaka2011}.  Under trivial typing the present
framework reproduces that semantic class; finite typing refines which observed
tuples may be identified while preserving compositional rule witnesses.
The MCFG proofs additionally require componentwise type propagation and, for
the quantitative boundary theorem, simultaneous control of several tuple
components.  Yoshinaka's $(k,\ell)$ hierarchy is a separate CFG result
\cite{Yoshinaka2008}; $h_\Sigma^{k,\ell}$ lifts the same boundary
information componentwise.  Since every $(f,h)$ condition includes arity one,
\[
  \Cmcf{f}{h}\cap\mathrm{CFL}\subseteq\Cmcf{1}{h}.
\]
The examples below separate finite typing from the trivial baseline,
prefix--suffix typings, and regular-filter explanations.

\subsection{Yoshinaka's Baseline and Prefix--Suffix Specialization}

\begin{definition}[Yoshinaka multidimensional substitutability]
\label{def:yoshinaka-pd-substitutability}
Using the sentence-context sets $\mathcal E_d$ defined in
Section~\ref{sec:prelim}, a language $L$ is \emph{$f$D-substitutable} if,
for every $1\le d\le f$, all $E,F\in\mathcal E_d$, and all
$\tuple x,\tuple y\in(\Sigma^+)^d$,
\[
 E[\tuple x],\quad E[\tuple y],\quad F[\tuple x]\in L
 \Longrightarrow
 F[\tuple y]\in L.
\]
Write $\YSub{f}$ for this semantic class and
$\YSL{f}{r}:=\YSub{f}\cap\mathsf{MCFL}(f,r)$, where $r$ bounds rule rank.
We refer to the displayed four-instance implication as the
\emph{substitution implication}.
Yoshinaka's original definition quantifies over every arity $d\le f$, uses
nonempty tuple components, and requires nonempty internal separators in
$d$-contexts; footnote~3 separately discusses the stronger alternative that
permits empty components and separators
\cite[Section~3.1 and footnote~3]{Yoshinaka2011}.
\end{definition}

Recall the finite monoid morphism $h_\Sigma^{k,\ell}$ from
Definition~\ref{def:prefix-suffix-typing}.  In Yoshinaka's word-level
definition, $w_1,y_1,w_1',w_2,y_2,w_2'\in\Sigma^*$,
$u\in\Sigma^k$, and $v\in\Sigma^\ell$, with
$uy_1v,uy_2v\ne\eps$; the condition is
\[
 w_1uy_1vw_1',\ w_1uy_2vw_1'\in L
 \quad\Longrightarrow\quad
 \bigl(w_2uy_1vw_2'\in L\iff w_2uy_2vw_2'\in L\bigr)
\]
\cite{Yoshinaka2008}.  Thus the compared fragments $uy_iv$ are nonempty even
though the middle words $y_i$ themselves may be empty.  Equality of their
$h_\Sigma^{k,\ell}$-values is exactly equality of the boundary information
used in this condition, including the exact identity of words shorter than
$m_0=\max\{1,k+\ell\}$.

\begin{proposition}[Semantic specializations of finite typing]
\label{prop:semantic-specializations}
The following hold.
\begin{enumerate}[label=(\roman*),leftmargin=*]
\item For the one-element typing $h_{\mathrm{triv}}$ defined above and every
$f\ge1$,
\[
  \Cmcf{f}{h_{\mathrm{triv}}}=\YSub{f}\cap f\text{-}\mathrm{MCFL},
\]
and for every finite typing $h$,
\[
  \YSub{f}\cap f\text{-}\mathrm{MCFL}\subseteq\Cmcf{f}{h}.
\]
\item At $f=1$, the prefix--suffix typing recovers Yoshinaka's word-level
$(k,\ell)$ condition: for every context-free language $L$,
\[
  L\text{ is $(k,\ell)$-substitutable}
  \quad\Longleftrightarrow\quad
  L\in\Cmcf{1}{h_\Sigma^{k,\ell}}.
\]
\end{enumerate}
\end{proposition}

\begin{proof}
For (i), trivial typing makes the type premise vacuous, and the inclusion for
general $h$ follows from $h_{\mathrm{triv}}\preceq h$.

For (ii), let nonempty $x,y$ have the same $h_\Sigma^{k,\ell}$-type.  If they
are shorter than $m_0$, then $x=y$; otherwise write
$x=uy_1v$ and $y=uy_2v$ with the common length-$k$ prefix $u$ and
length-$\ell$ suffix $v$.  Yoshinaka's $(k,\ell)$ condition then gives
$\D_L(x)=\D_L(y)$.  Conversely, in the premises of that condition the two
nonempty compared fragments have the same prefix--suffix type, so
$(1,h_\Sigma^{k,\ell})$-tuple substitutability transfers acceptance between
the two outer contexts in both directions.
\end{proof}

For general $f$, componentwise $h_\Sigma^{k,\ell}$-equality is exactly the
corresponding multidimensional rectangular-window condition.

\paragraph{Compatibility with Yoshinaka's bounded-rank scale.}
For trivial typing every productive nonterminal has a unique componentwise
type, so Corollary~\ref{cor:no-splitting-thickness}, rather than the more
general boundary theorem, gives the sharper comparison.  For a reduced good
rank-$\le r$ grammar,
\[
 |\Wit(\widetilde G_0)|\le |V|+2|P|,
 \qquad
 \|\Wit(\widetilde G_0)\|_+
 \le (|V|+2|P|)\bigl(1+(r+2)\max\{1,\theta_G\}\bigr).
\]
Yoshinaka's characteristic set satisfies
$|K_G|\le |P|$ and $\|K_G\|\le |P|\tau_G$
\cite[Lemma~7]{Yoshinaka2011}, and his Theorem~1 gives polynomial time and
data for fixed dimension and rank~\cite[Theorem~1]{Yoshinaka2011}.  Thus the
present trivial-typing specialization recovers the same grammar-size--rule-
thickness efficiency scale, while the exact constants and witness
organization differ.

\paragraph{Same-typing regular-filter fragment.}
For a fixed finite typing $h:\Sigma^*\to M$, define
\[
  \Ffrag{f}{h}
  :=\bigl\{L_0\cap h^{-1}(\mathsf{Acc}):L_0\in\YSub{f},\ \mathsf{Acc}\subseteq M\bigr\}
    \cap f\text{-}\mathrm{MCFL}.
\]
This is the part of the typed class obtained by taking Yoshinaka's trivial-typing
baseline and applying only a regular filter recognized by the \emph{same}
morphism $h$.

\begin{proposition}[Same-typing filtering is sufficient]
\label{prop:same-typing-filter-fragment}
For every fixed $f$ and finite typing $h$,
\[
  \Ffrag{f}{h}\subseteq\Cmcf{f}{h}.
\]
\end{proposition}

\begin{proof}
For $L=L_0\cap h^{-1}(\mathsf{Acc})$, $f$D-substitutability transfers the
$L_0$ part between equal-$h$-type tuples, while compositionality of $h$
preserves the whole-word filter value in every fixed context.  Symmetry gives
equal tuple distributions.
\end{proof}

The same-typing qualifier excludes arbitrary target-dependent regular filters;
these are treated after Theorem~\ref{thm:every-fanout-strictness}.


\subsection{An Intrinsically Rank-Two Boundary-Typed Target}
\label{subsec:intrinsic-ranktwo}

The next example makes the distinction between the rank-one and higher-rank
quantitative regimes intrinsic to the language.  Its rule-rank obstruction is
not created by finite typing: the balanced-count factor used below already
belongs to Yoshinaka's trivial-typing class.  The construction instead combines
that obstruction with minimum fan-out two and a separation from the
trivial-typing class inside one boundary-typed target.

Let
\[
 \Sigma_{\mathrm{intr}}
 :=\{\mathtt a,\mathtt b,\mathtt l,\mathtt r,\mathtt p,\mathtt q,
       \mathtt c,\mathtt d,\mathtt{\#},\mathtt{\$},
       \mathtt u,\mathtt v,\mathtt w\}
\]
and put
\[
 P_{\mathrm{intr}}
 :=\{\mathtt{lpc},\mathtt{lpd},\mathtt{lqc},
       \mathtt{rpc},\mathtt{rpd}\}.
\]
Define
\[
 \mathrm{MIX}_2^+
 :=\{x\in\{\mathtt a,\mathtt b\}^+:
       |x|_{\mathtt a}=|x|_{\mathtt b}\},
 \qquad
 T_3^+
 :=\{\mathtt u^n\mathtt v^n\mathtt w^n:n\ge1\},
\]
and
\[
 L_{\mathrm{intr}}
 :=P_{\mathrm{intr}}\,\mathtt{\#}\,\mathrm{MIX}_2^+\,
   \mathtt{\$}\,T_3^+.
\]

To make the role of the witness explicit, let
\[
 \varphi_{\mathrm{mix}}:\{\mathtt a,\mathtt b\}^*\to\mathbb Z,
 \qquad
 \varphi_{\mathrm{mix}}(\mathtt a)=1,\quad
 \varphi_{\mathrm{mix}}(\mathtt b)=-1.
\]
Then
\[
 \mathrm{MIX}_2^+
 =\{\mathtt a,\mathtt b\}^+\cap\varphi_{\mathrm{mix}}^{-1}(0).
\]
The language $\{\mathtt a,\mathtt b\}^+$ is trivially
$f$-tuple-substitutable, so
Proposition~\ref{prop:abelian-coset-filter}(ii) gives
$\mathrm{MIX}_2^+\in\YSub{f}$ for every $f\ge1$; the rank-two CFG for $B$
displayed below also shows that this ingredient is context-free.  Thus the
intrinsic rank-two phenomenon used here already occurs inside the
trivial-typing baseline.  The role of $L_{\mathrm{intr}}$ is to combine that
known rule-rank obstruction with non-context-freeness, minimum fan-out two,
and a separation from $\YSub{f}$ in a single finitely typed target.

\begin{theorem}[A non-context-free intrinsic-rank-two typed witness]
\label{thm:intrinsic-ranktwo-witness}
The language $L_{\mathrm{intr}}$ satisfies
\[
 L_{\mathrm{intr}}\in\mathsf{MCFL}(2,2)\setminus\mathrm{CFL},
\]
and, for every $f\ge2$,
\[
 L_{\mathrm{intr}}
 \in\Cmcf{f}{h_{\Sigma_{\mathrm{intr}}}^{1,1}}
 \setminus\YSub{f}.
\]
Moreover, for every finite $f'\ge1$,
\[
 L_{\mathrm{intr}}\notin\mathsf{MCFL}(f',1).
\]
Consequently the minimum fan-out of $L_{\mathrm{intr}}$ is exactly two, and
its minimum rule rank among all finite-fan-out MCFG presentations is exactly
two.  In particular, Theorem~\ref{thm:boundary-polydata} applies to a
non-context-free language that cannot be covered by
Theorem~\ref{thm:rankone-polydata} at any finite fan-out.
\end{theorem}

\begin{proof}
First we give a reduced good fan-out-two grammar of maximum rule rank two.
Let $B$ generate the nonempty balanced-count language by
\[
\begin{aligned}
 B\to{}&
 \mathtt a B\mathtt b B
 \mid \mathtt b B\mathtt a B
 \mid \mathtt a B\mathtt b
 \mid \mathtt a\mathtt b B
 \mid \mathtt a\mathtt b\\
 &{}\mid
 \mathtt b B\mathtt a
 \mid \mathtt b\mathtt a B
 \mid \mathtt b\mathtt a .
\end{aligned}
\]
The usual first-return induction on
$|x|_{\mathtt a}-|x|_{\mathtt b}$ gives
$L_B=\mathrm{MIX}_2^+$.

Let the fan-out-two nonterminal $A$ have rules
\[
 A\to(\mathtt u,\mathtt w)
 \mid
 (\mathtt u\xi_1^1\mathtt v,\mathtt w\xi_1^2)(A).
\]
It derives exactly the tuples
\[
 \{(\mathtt u^n\mathtt v^{n-1},\mathtt w^n):n\ge1\}.
\]
For every $\alpha\in P_{\mathrm{intr}}$, add the start rule
\[
 S\to
 \alpha\mathtt{\#}\xi_1^1\mathtt{\$}
 \xi_2^1\mathtt v\xi_2^2(B,A).
\]
The resulting grammar generates exactly $L_{\mathrm{intr}}$.  It is reduced,
$\lambda$-free, linear, nondeleting, and nonpermuting.  It is also nonmerging:
the two components of $A$ are separated by the literal terminal $\mathtt v$
in every start rule.  Hence it is a good fan-out-two grammar of maximum rule
rank two.

The language is not context-free.  Intersect it with the regular language
\[
 R:=
 \mathtt{lpc\#ab\$}\,
 \mathtt u^+\mathtt v^+\mathtt w^+ .
\]
Then
\[
 L_{\mathrm{intr}}\cap R
 =
 \{\mathtt{lpc\#ab\$}\mathtt u^n\mathtt v^n\mathtt w^n:n\ge1\}.
\]
A homomorphism erasing the fixed prefix and fixing
$\mathtt u,\mathtt v,\mathtt w$ maps this language onto
$\{\mathtt u^n\mathtt v^n\mathtt w^n:n\ge1\}$, which is not context-free.
Closure of context-free languages under intersection with regular languages
and homomorphism therefore shows that $L_{\mathrm{intr}}$ is not
context-free.  Since fan-out-one MCFGs are exactly context-free grammars in
the present syntax, while the displayed grammar has fan-out two, the minimum
fan-out is exactly two.

For typed membership, let
\[
 Q_{\mathrm{intr}}
 :=P_{\mathrm{intr}}\,\mathtt{\#}\,\{\mathtt a,\mathtt b\}^+\,
   \mathtt{\$}\,\mathtt u^+\mathtt v^+\mathtt w^+ .
\]
This language is strictly $2$-local.  One complete set of allowed adjacent
pairs, including end markers, is
\[
\begin{gathered}
 \vdash\mathtt l,\ \vdash\mathtt r,\quad
 \mathtt{lp},\ \mathtt{lq},\ \mathtt{rp},\quad
 \mathtt{pc},\ \mathtt{pd},\ \mathtt{qc},\\
 \mathtt{c\#},\ \mathtt{d\#},\quad
 \mathtt{\#a},\ \mathtt{\#b},\quad
 \mathtt{aa},\ \mathtt{ab},\ \mathtt{ba},\ \mathtt{bb},\\
 \mathtt{a\$},\ \mathtt{b\$},\quad
 \mathtt{\$u},\quad
 \mathtt{uu},\ \mathtt{uv},\ \mathtt{vv},\ \mathtt{vw},\
 \mathtt{ww},\quad
 \mathtt w\dashv .
\end{gathered}
\]
Define the homomorphism
\[
 \Phi_{\mathrm{intr}}:\Sigma_{\mathrm{intr}}^*\to\mathbb Z^3
\]
by
\[
\begin{aligned}
 \Phi_{\mathrm{intr}}(\mathtt a)&=(1,0,0),&
 \Phi_{\mathrm{intr}}(\mathtt b)&=(-1,0,0),\\
 \Phi_{\mathrm{intr}}(\mathtt u)&=(0,1,0),&
 \Phi_{\mathrm{intr}}(\mathtt v)&=(0,-1,1),&
 \Phi_{\mathrm{intr}}(\mathtt w)&=(0,0,-1),
\end{aligned}
\]
and let every other terminal have image $(0,0,0)$.  Then
\[
 L_{\mathrm{intr}}
 =Q_{\mathrm{intr}}\cap\Phi_{\mathrm{intr}}^{-1}(0).
\]
Indeed, the three coordinates impose respectively
\[
 |x|_{\mathtt a}=|x|_{\mathtt b},\qquad
 |x|_{\mathtt u}=|x|_{\mathtt v},\qquad
 |x|_{\mathtt v}=|x|_{\mathtt w}.
\]
Lemma~\ref{lem:sl2-boundary-control} and
Proposition~\ref{prop:abelian-coset-filter}(ii) show that the displayed
intersection is $(f,h_{\Sigma_{\mathrm{intr}}}^{1,1})$-tuple-substitutable
for every $f\ge1$.  Together with the fan-out-two grammar above, this gives
\[
 L_{\mathrm{intr}}
 \in\Cmcf{f}{h_{\Sigma_{\mathrm{intr}}}^{1,1}}
 \qquad\text{for every }f\ge2.
\]

The language is outside Yoshinaka's trivial-typing class already at arity one.
Take
\[
 x:=\mathtt p,\qquad y:=\mathtt q,
\]
and one-hole contexts
\[
 E:=\mathtt l\blank_1\mathtt{c\#ab\$uvw},
 \qquad
 F:=\mathtt r\blank_1\mathtt{d\#ab\$uvw}.
\]
Then
\[
 E[x],\ E[y],\ F[x]\in L_{\mathrm{intr}},
 \qquad
 F[y]\notin L_{\mathrm{intr}}.
\]
Thus the arity-one substitution implication fails, so
$L_{\mathrm{intr}}\notin\YSub{f}$ for every $f\ge1$.  Notice again that the
boundary typing distinguishes the offending one-letter factors:
\[
 h_{\Sigma_{\mathrm{intr}}}^{1,1}(\mathtt p)
 \ne h_{\Sigma_{\mathrm{intr}}}^{1,1}(\mathtt q).
\]

It remains to prove intrinsic rule rank.  Put
\[
 \mathrm{MIX}_2
 :=\{x\in\{\mathtt a,\mathtt b\}^*:
       |x|_{\mathtt a}=|x|_{\mathtt b}\}.
\]
Gallot calls an MCFG \emph{non-branching} when every positive-rank rule has a
single child nonterminal; in the present terminology this is exactly rule
rank at most one, with no globally fixed fan-out bound
\cite[Definition~38]{Gallot2021}.  Since every finite grammar has a finite
maximum nonterminal arity, Gallot's non-branching MCFL class is
\[
 \bigcup_{f'<\infty}\mathsf{MCFL}(f',1).
\]
Gallot's Theorem~28 allows a non-branching grammar to be taken non-erasing and
non-permuting without changing the generated class.  Proposition~5 then shows
that if $\mathrm{MIX}_2$ were generated by a non-branching MCFG, all of its
words would admit a uniformly bounded $k$-derivation, whereas Proposition~6
constructs, for every $k$, a word with no such derivation.  Thus
$\mathrm{MIX}_2$ is not non-branching; this is the MCFG lower bound used in
his Theorem~31, with Theorem~29 identifying non-branching MCFLs with
finite-index EDT0L languages
\cite[Theorems~28--29, Propositions~5--6, Theorem~31]{Gallot2021}.
Bishop et al. later give a peer-reviewed treatment of the same obstruction
in the closely corresponding R-MCFG formalism: Definition~6.1 gives the
rank-zero/unary rule shape, Lemma~6.8 supplies the normal form, and Section~7
rules out the same balanced language.  Their Theorem~C strengthens the
conclusion to non-EDT0L
\cite[Definition~6.1, Lemma~6.8, Section~7, Theorem~C]{BishopEtAl2026EDT0L}.
Together with the rank-preserving good normalization of
Proposition~\ref{prop:good-normalization}, this gives an independent
corroborating route; Gallot's direct non-branching MCFG lower bound is the
implication used here.

Now let $\pi$ fix $\mathtt a,\mathtt b$ and erase every other terminal in
$\Sigma_{\mathrm{intr}}$.  Then
\[
 \pi(L_{\mathrm{intr}})=\mathrm{MIX}_2^+.
\]
If $L_{\mathrm{intr}}$ had a fan-out-$f'$ rule-rank-one presentation,
applying $\pi$ to the terminal material in its rule templates would preserve
fan-out and rule rank and would give such a presentation of
$\mathrm{MIX}_2^+$.  Empty terminal pieces created by erasure are permitted
in the standard MCFG formalism.  Adding a fresh start symbol with one unary
identity link to the old start and one rank-zero rule generating $\eps$ would
then give a fan-out-$f'$ rule-rank-one presentation of $\mathrm{MIX}_2$,
contradicting Gallot's non-branching lower bound.  Hence
\[
 L_{\mathrm{intr}}\notin\mathsf{MCFL}(f',1)
 \qquad\text{for every finite }f'\ge1.
\]
Together with the displayed rank-two grammar, this proves the claimed
intrinsic rule rank.
\end{proof}

The construction is deliberately modular: the $\mathrm{MIX}_2^+$ factor
carries the obstruction to rule rank one, while the $T_3^+$ factor forces the
target outside CFL and hence to minimum fan-out two.  Thus the theorem
establishes the coexistence of the two lower bounds in a single finitely typed
target; the two bounds arise from separate factors rather than from an
irreducible interaction between the resources.

The non-branching lower bound itself is prior work.  What is specific to the
present framework is that a strictly local wrapper and the $(1,1)$ boundary
typing place this target inside the typed class, while the displayed
substitution rectangle keeps it outside Yoshinaka's trivial-typing baseline.
The wrapper retains the prior non-branching obstruction, so the same target
also has intrinsic rule rank two across all finite fan-outs.

\subsection{Finite Typings beyond the Trivial Baseline}
\label{subsec:finite-typing-separations}


The filtering results of Section~\ref{sec:prelim} give useful sufficient
constructions, but finite typing is not merely regular post-filtering of the
trivial-typing class.  The next linear example gives a stronger separation and,
at the same time, a sharp fan-out effect.

Let
\[
  L_{\mathrm{cd}}:=
  \{\mathtt a^n\mathtt c\mathtt b^n:n\ge0\}
  \cup
  \{\mathtt a^n\mathtt d\mathtt b^{2n}:n\ge0\}
\]
over $\Sigma_{\mathrm{cd}}:=\{\mathtt a,\mathtt b,\mathtt c,\mathtt d\}$.  Let
$M_{\mathrm{cd}}:=\{1,\eta_{\mathtt c},\eta_{\mathtt d},0\}$, where $1$ is the identity, $0$ is a zero,
and every product of two elements of $\{\eta_{\mathtt c},\eta_{\mathtt d},0\}$ is $0$.  Define
\[
 h_{\mathrm{cd}}(\mathtt a)=h_{\mathrm{cd}}(\mathtt b)=1,
 \qquad
 h_{\mathrm{cd}}(\mathtt c)=\eta_{\mathtt c},
 \qquad
 h_{\mathrm{cd}}(\mathtt d)=\eta_{\mathtt d}.
\]

\begin{proposition}[A fan-out-sensitive typed language]
\label{prop:strong-fanout-separation}
The language $L_{\mathrm{cd}}$ is linear and
\[
  L_{\mathrm{cd}}\in\Cmcf{1}{h_{\mathrm{cd}}}.
\]
For every $k,\ell\ge0$,
\[
  L_{\mathrm{cd}}\notin\Cmcf{1}{h_{\Sigma_{\mathrm{cd}}}^{k,\ell}}.
\]
There do not exist a language $L_0\in \YSub{1}$ and a regular language
$L_{\mathrm{reg}}$ such that
$L_{\mathrm{cd}}=L_0\cap L_{\mathrm{reg}}$.  Moreover, for every finite monoid
homomorphism $g:\Sigma_{\mathrm{cd}}^*\to N$,
\[
  L_{\mathrm{cd}}\notin\Cmcf{2}{g}.
\]
\end{proposition}

\begin{proof}
The linear grammar
\[
  S\to C\mid D,\qquad
  C\to\mathtt c\mid\mathtt a C\mathtt b,\qquad
  D\to\mathtt d\mid\mathtt a D\mathtt b\mathtt b
\]
generates $L_{\mathrm{cd}}$.

For $(1,h_{\mathrm{cd}})$-substitutability, let equal-type nonempty factors
$x,y$ share an accepting context.  Type $1$ contains no center symbol, so the
context fixes the branch and side and the count constraint forces $x=y$.
For type $\eta_{\mathtt c}$, writing
$x=\mathtt a^{n_1}\mathtt c\mathtt b^{n_2}$ and
$y=\mathtt a^{n_1'}\mathtt c\mathtt b^{n_2'}$, common acceptance gives
$n_1-n_2=n_1'-n_2'$, the invariant for every other $\mathtt c$-context.
For type $\eta_{\mathtt d}$ the invariant is
$2n_1-n_2=2n_1'-n_2'$, while type $0$ has empty distribution.  Hence
$L_{\mathrm{cd}}\in\Cmcf{1}{h_{\mathrm{cd}}}$.

Fix $k,\ell$ and choose $m\ge\max\{1,k+\ell\}$.  The words
$x=\mathtt a^m\mathtt c\mathtt b^m$ and
$y=\mathtt a^m\mathtt d\mathtt b^{2m}$ have the same
$h_{\Sigma_{\mathrm{cd}}}^{k,\ell}$-type and share the empty context, but
$\mathtt a\blank_1\mathtt b$ accepts $x$ and rejects $y$.  Thus every
prefix--suffix typing fails at fan-out one.

Suppose $L_{\mathrm{cd}}=L_0\cap L_{\mathrm{reg}}$ with
$L_0\in\YSub{1}$ and $L_{\mathrm{reg}}$ regular.  Since
$\mathtt c,\mathtt d\in L_0$, substitutability gives
$\D_{L_0}(\mathtt c)=\D_{L_0}(\mathtt d)$; using
$\mathtt a\mathtt d\mathtt b^2,\mathtt a\mathtt c\mathtt b\in L_0$ then gives
$\D_{L_0}(\mathtt c\mathtt b)=\D_{L_0}(\mathtt c\mathtt b^2)$.
Consequently
\[
 \mathtt a^m\mathtt c\mathtt b^{m+t}\in L_0
 \qquad(m\ge1,\ t\ge0).
\]
For $m$ larger than the state number of a DFA for $L_{\mathrm{reg}}$, pumping
the final $\mathtt b^m$ block of the accepted target word gives some $t>0$
with the same word in $L_{\mathrm{reg}}$, contradicting the assumed
intersection representation.

Finally, for finite $g:\Sigma_{\mathrm{cd}}^*\to N$, choose $1\le i<j$ with
\[
 g(\mathtt a^i)=g(\mathtt a^j),\qquad
 g(\mathtt b^i)=g(\mathtt b^j).
\]
The tuples $(\mathtt a^i,\mathtt b^i)$ and
$(\mathtt a^j,\mathtt b^j)$ have equal componentwise $g$-type and are both
accepted by $\blank_1\mathtt c\blank_2$, whereas
$\blank_1\mathtt d\blank_2\mathtt b^i$ accepts only the first:
$\mathtt a^j\mathtt d\mathtt b^{i+j}\notin L_{\mathrm{cd}}$ for $i<j$.
Thus arity two fails for every finite $g$.
\end{proof}

Hence $L_{\mathrm{cd}}$ separates finite typing from every prefix--suffix
typing and arbitrary regular filtering at fan-out one, while no finite typing
repairs the language once arity two is tested.

\begin{proposition}[The copy language has no finite typing]
\label{prop:copy-infinite-typing}
Let $\Sigma:=\{\mathtt a,\mathtt b\}$ and
$\mathrm{COPY}:=\{ww:w\in\Sigma^+\}$.  Then
$\mathrm{COPY}\in2\text{-}\mathrm{MCFL}$, but for every finite monoid
homomorphism $h:\Sigma^*\to M$ and every $f\ge1$, COPY is not
$(f,h)$-tuple-substitutable.  In fact the failure already occurs at arity one.
\end{proposition}

\begin{proof}
The standard fan-out-two grammar
\[
\begin{aligned}
 A&\to(\mathtt a,\mathtt a)\mid(\mathtt b,\mathtt b)
      \mid(\mathtt a \xi_1^1,\mathtt a \xi_1^2)(A)
      \mid(\mathtt b \xi_1^1,\mathtt b \xi_1^2)(A),\\
 S&\to(\xi_1^1\xi_1^2)(A)
\end{aligned}
\]
generates COPY, so $\mathrm{COPY}\in2\text{-}\mathrm{MCFL}$.

Fix finite $h:\Sigma^*\to M$ and choose $n$ with $2^n>|M|$ and distinct
$u,v\in\Sigma^n$ satisfying $h(u)=h(v)$.  Then $x:=uu$ and $y:=vv$ have equal
$h$-type and are both accepted by the empty context.  The context
$E:=u\blank_1u$ accepts $x$, since $E[x]=uuuu$.  If it accepted $y$, then
$uvvu\in\mathrm{COPY}$ would have equal halves, so $uv=vu$; because
$|u|=|v|$, this forces $u=v$, a contradiction.  Thus arity-one
substitutability fails for every finite $h$, and hence so does the $(f,h)$
condition for every $f\ge1$.
\end{proof}

\subsection{Strictness and Typings beyond Prefix--Suffix Windows}
\label{subsec:comparison-separations}

The prefix--suffix typings do not exhaust the finite-typing framework.
Let $\Sigma_Q:=\{\mathtt 0,\mathtt 1,\mathtt\#,\mathtt a,\mathtt b\}$ and
\[
  Q:=\{w\mathtt\#\mathtt a^i\mathtt b^j:w\in\{\mathtt 0,\mathtt 1\}^*,\ i,j\ge0,\ \text{$w$ contains $\mathtt{11}$}\}.
\]
Fix a complete DFA for $Q$, let $h_Q:\Sigma_Q^*\to M_Q$ be its transition
morphism, and let $\mathsf{Acc}_Q\subseteq M_Q$ satisfy $Q=h_Q^{-1}(\mathsf{Acc}_Q)$.

\begin{proposition}[A non-boundary typing beyond every prefix--suffix window]
\label{prop:nonboundary-rankone-separator}
For every $s\ge1$, let
\[
 J_s:=\{w\mathtt\#\mathtt a^{sn}\mathtt b^n:
       w\in\{\mathtt 0,\mathtt 1\}^n,
       \text{$w$ contains $\mathtt{11}$}\}.
\]
Then $J_s$ is a non-context-free fan-out-two rank-one language satisfying
\[
  J_s\in\Cmcf{2}{h_Q},
  \qquad
  J_s\notin\Cmcf{2}{h_{\Sigma_Q}^{k,\ell}}
  \quad\text{for every }k,\ell\ge0.
\]
\end{proposition}

\begin{proof}
A fan-out-two rank-one grammar is
\[
\begin{aligned}
 A&\to(\mathtt{11},\mathtt a^{2s}\mathtt b^2),\\
 A&\to(\mathtt 0\xi_1^1,\mathtt a^s\xi_1^2\mathtt b)(A)
   \mid(\mathtt 1\xi_1^1,\mathtt a^s\xi_1^2\mathtt b)(A),\\
 A&\to(\xi_1^1\mathtt 0,\mathtt a^s\xi_1^2\mathtt b)(A)
   \mid(\xi_1^1\mathtt 1,\mathtt a^s\xi_1^2\mathtt b)(A),\\
 S&\to(\xi_1^1\mathtt\#\xi_1^2)(A).
\end{aligned}
\]
It generates $J_s$.  Intersecting with
$\mathtt 1^*\mathtt\#\mathtt a^*\mathtt b^*$ gives
$\{\mathtt 1^n\mathtt\#\mathtt a^{sn}\mathtt b^n:n\ge2\}$, which is
not context-free: otherwise the indicated block homomorphism, inverse image,
and regular intersection would yield
$\{\mathtt x^n\mathtt\#\mathtt y^n\mathtt z^n:n\ge2\}$.

Moreover $J_s=Q\cap\varphi_s^{-1}(0,0)$ for the abelian map
$\varphi_s(\mathtt 0)=\varphi_s(\mathtt 1)=(-s,-1)$,
$\varphi_s(\mathtt a)=(1,0)$, $\varphi_s(\mathtt b)=(0,1)$, and
$\varphi_s(\mathtt\#)=(0,0)$, so
Proposition~\ref{prop:abelian-coset-filter} gives
$J_s\in\Cmcf{2}{h_Q}$.

Fix $k,\ell\ge0$, put $n_0:=k+\ell+4$, and take
$x=\mathtt 0^k\mathtt{0110}\mathtt 0^\ell$ and
$y=\mathtt 0^k\mathtt{0100}\mathtt 0^\ell$.  With
\[
 E=\mathtt{11}\blank_1\mathtt\#\mathtt a^{s(n_0+2)}\mathtt b^{n_0+2},
 \qquad F=\blank_1\mathtt\#\mathtt a^{sn_0}\mathtt b^{n_0},
\]
the factors $x,y$ have the same $h_{\Sigma_Q}^{k,\ell}$-type and
$E[x],E[y],F[x]\in J_s$, whereas $F[y]\notin J_s$.  Thus the
$(2,h_{\Sigma_Q}^{k,\ell})$ condition already fails at arity one.
\end{proof}

Thus the arbitrary-$h$ rank-one theorem reaches beyond every fixed
prefix--suffix typing.  Since
$J_s=\varphi_s^{-1}(0,0)\cap h_Q^{-1}(\mathsf{Acc}_Q)$, this witness concerns
non-boundary typing, not regular-filter separation.  We next obtain strictness
at every fan-out with a separator-free witness; Yoshinaka's Example~3 uses
separator-marked equal-block languages instead~\cite[Example~3]{Yoshinaka2011}.

Following the standard local-testability terminology~\cite{McNaughtonPapert1971}, call
$L_{\mathrm{sl}}\subseteq\Sigma^*$ \emph{strictly 2-local} if there is a set
$\mathsf{Adj}\subseteq(\Sigma\cup\{\vdash,\dashv\})^2$ such that
\[
 w\in L_{\mathrm{sl}}
 \quad\Longleftrightarrow\quad
 \text{every adjacent pair in $\vdash w\dashv$ belongs to $\mathsf{Adj}$},
\]
where $\vdash,\dashv\notin\Sigma$ are fixed boundary markers.

\begin{lemma}[Strictly 2-local constraints are boundary typed]
\label{lem:sl2-boundary-control}
If $L_{\mathrm{sl}}\subseteq\Sigma^*$ is strictly 2-local, then for every $f\ge1$,
\[
  L_{\mathrm{sl}}\in\Cmcf{f}{h_\Sigma^{1,1}}.
\]
\end{lemma}

\begin{proof}
Let equal-$h_\Sigma^{1,1}$-type nonempty tuples $\tuple x,\tuple y$ share an
accepting context $E_0$.  Internal adjacent pairs of each $y_i$ are permitted
because they occur in $E_0[\tuple y]$.  If $E[\tuple x]\in L_{\mathrm{sl}}$,
every adjacent pair of $E[\tuple y]$ is either in fixed material, internal to
a $y_i$, or crosses a component boundary.  These are permitted respectively
by $E[\tuple x]$, by $E_0[\tuple y]$, or by equality of the first/last letters
encoded by the boundary type; end-marker pairs are unchanged.  Thus
$E[\tuple y]\in L_{\mathrm{sl}}$, and symmetry finishes the proof.
\end{proof}

\begin{theorem}[Strictness at every fan-out]
\label{thm:every-fanout-strictness}
For $m\ge2$ let
$\Sigma_m:=\{\mathtt a_1,\ldots,\mathtt a_{2m}\}$ and
\[
  H_m:=\{\mathtt a_1^n\mathtt a_2^n\cdots\mathtt a_{2m}^n:n\ge1\}.
\]
Then $H_m\in\mathsf{MCFL}(m,1)$ but $H_m$ is not an $(m-1)$-MCFL.  Moreover,
for every $f\ge m$,
\[
  H_m\in\Cmcf{f}{h_{\Sigma_m}^{1,1}}
  \setminus\bigl(\YSub{f}\cup\Ffrag{f}{h_{\Sigma_m}^{1,1}}\bigr).
\]
\end{theorem}

The alphabet $\Sigma_m$ varies with $m$; the hierarchy content of the theorem
is that the witness has minimum fan-out exactly $m$.  For strictness at any
single fixed fan-out $f\ge2$, the case $H_2$ already supplies a witness for
every tested fan-out at least two.

\begin{proof}
A good rank-one grammar is
\[
\begin{aligned}
  A&\to(\mathtt a_1,\mathtt a_3,\ldots,\mathtt a_{2m-1}),\\
  A&\to(\mathtt a_1\xi_1^1\mathtt a_2,
        \mathtt a_3\xi_1^2\mathtt a_4,\ldots,
        \mathtt a_{2m-1}\xi_1^m\mathtt a_{2m})(A),\\
  S&\to(\xi_1^1\mathtt a_2\,\xi_1^2\mathtt a_4\cdots
         \xi_1^m\mathtt a_{2m})(A).
\end{aligned}
\]
It derives from $A$ the tuples
$(\mathtt a_1^n\mathtt a_2^{n-1},\ldots,
  \mathtt a_{2m-1}^n\mathtt a_{2m}^{n-1})$ and hence generates $H_m$.
Seki et al.'s Lemma~3.3, applied with their parameter $m-1$,
shows that
$\{\mathtt a_1^n\cdots\mathtt a_{2m-1}^n:n\ge1\}$ is not an
$(m-1)$-MCFL~\cite[Lemma~3.3, p.~202]{SekiEtAl1991}.  Since MCFLs are closed under homomorphism as a consequence of their
substitution closure~\cite[Theorem~3.9, p.~203]{SekiEtAl1991}, erasing
$\mathtt a_{2m}$ proves that $H_m$ is not an $(m-1)$-MCFL.

For membership, put
$Q_m:=\mathtt a_1^*\cdots\mathtt a_{2m}^*\cap\Sigma_m^+$.
This language is strictly $2$-local, and
\[
  H_m=Q_m\cap\varphi_m^{-1}(0),\qquad
  \varphi_m(w):=(|w|_{\mathtt a_i}-|w|_{\mathtt a_{i+1}})_{1\le i<2m}.
\]
Lemma~\ref{lem:sl2-boundary-control} and
Proposition~\ref{prop:abelian-coset-filter}(ii) therefore give
$H_m\in\Cmcf{f}{h_{\Sigma_m}^{1,1}}$ for every $f\ge m$.

It remains to exclude Yoshinaka's substitution implication.  Put
$d:=\lceil m/2\rceil$.  Since $d\le2m-2d\le2d$, choose
$\ell_i\in\{1,2\}$ with $\sum_i\ell_i=2m-2d$, and set
\[
\begin{aligned}
 q_i&:=2+2(i-1)+\sum_{j<i}\ell_j,&
 x_i&:=\mathtt a_{q_i}\cdots\mathtt a_{q_i+\ell_i-1},\\
 v_0&:=\mathtt a_1,&v_d&:=\mathtt a_{2m},&
 v_i&:=\mathtt a_{q_i+\ell_i}\mathtt a_{q_i+\ell_i+1}\quad(1\le i<d).
\end{aligned}
\]
Then $w_1:=\mathtt a_1\cdots\mathtt a_{2m}=v_0x_1v_1\cdots x_dv_d$.
With $E:=v_0\blank_1v_1\cdots\blank_dv_d$, we have $E[\tuple x]=w_1$.

In $w_2:=\mathtt a_1^2\cdots\mathtt a_{2m}^2$, place $v_0,v_d$ at the two
ends and each interior $v_i=\mathtt a_j\mathtt a_{j+1}$ across its block
boundary.  The intervening nonempty segments form $\tuple y$ with
$E[\tuple y]=w_2$ and $y_1$ beginning in $\mathtt a_1$.  Place each $x_i$ in
$w_2$ starting at the last occurrence of its first letter (straddling the
block boundary when $|x_i|=2$); the remaining material defines
$F\in\mathcal E_d$ with $F[\tuple x]=w_2$.  Its fixed prefix ends in the first
$\mathtt a_2$, so $F[\tuple y]$ contains $\mathtt a_2\mathtt a_1$ and is not
in $H_m$.  Hence
\[
 E[\tuple x],E[\tuple y],F[\tuple x]\in H_m,\qquad
 F[\tuple y]\notin H_m,
\]
with $d\le m\le f$, proving $H_m\notin\YSub{f}$.

The same rectangle excludes $\Ffrag{f}{h_{\Sigma_m}^{1,1}}$: all four complete
words have the same first and last symbols and length at least two, hence the
same whole-word boundary type.  The first three would force the fourth into
both $L_0$ and the same filter, a contradiction.
\end{proof}

Consequently, for every $m\ge2$, every $f\ge m$, and every $r\ge1$,
\[
  \YSL{f}{r}\subsetneq
  \Cmcf{f}{h_{\Sigma_m}^{1,1}}\cap\mathsf{MCFL}(f,r).
\]
The witness $H_m$ has minimum fan-out $m$ and rule rank one.

\paragraph{Arbitrary target-dependent regular filters.}
For comparison, on a fixed alphabet define
\[
 \mathcal F_f^{\mathrm{arb}}
 :=\bigl\{L_0\cap R:L_0\in\YSub{f},\ R\text{ regular}\bigr\}
   \cap f\text{-}\mathrm{MCFL}.
\]
Because every regular language has a finite transition morphism,
\[
 \mathcal F_f^{\mathrm{arb}}=\bigcup_h\Ffrag{f}{h}.
\]
This union contains all regular languages (take $\Sigma^*\in\YSub{f}$), so as
one positive-data class it is not TxtEx-identifiable by Gold's obstruction.

The witnesses $H_m$ and $L_{\mathrm{intr}}$ lie in this larger fragment:
\[
 H_m=\varphi_m^{-1}(0)\cap\mathtt a_1^+\cdots\mathtt a_{2m}^+,\qquad
 L_{\mathrm{intr}}=\Phi_{\mathrm{intr}}^{-1}(0)\cap Q_{\mathrm{intr}}.
\]
They therefore witness fan-out/rank and same-typing separation, not arbitrary
regular-filter separation.  Proposition~\ref{prop:strong-fanout-separation}
gives the latter at fan-out one; the next theorem keeps it inside a finitely
typed fan-out-two rank-one target.

\begin{theorem}[A fan-out-two typed language beyond arbitrary regular filtering]
\label{thm:arbitrary-filter-separation}
Let $\Sigma_\star:=\{\mathtt a,\mathtt b,\mathtt c,\mathtt d\}$ and
\[
 L_\star
 :=\{\mathtt a^n\mathtt c^{n+1}\mathtt b^n:n\ge0\}
   \cup
   \{\mathtt a^n\mathtt d^{n+1}\mathtt b^{2n}:n\ge0\}.
\]
Let
$g_\star:\Sigma_\star^*\to(\mathcal P(\{\mathtt c,\mathtt d\}),\cup)$
be the homomorphism defined by
\[
 g_\star(\mathtt a)=g_\star(\mathtt b)=\varnothing,
 \qquad
 g_\star(\mathtt c)=\{\mathtt c\},
 \qquad
 g_\star(\mathtt d)=\{\mathtt d\},
\]
and put
\[
 h_\star:=(h_{\Sigma_\star}^{1,1},g_\star).
\]
Then $L_\star$ is non-context-free, has minimum fan-out two and minimum
fan-out-two rule rank one, and for every $f\ge2$,
\[
  L_\star\in\Cmcf{f}{h_\star}
  \setminus\mathcal F_f^{\mathrm{arb}}.
\]
Equivalently, there are no $L_0\in\YSub{f}$ and regular $R_{\mathrm{reg}}$ such that
$L_\star=L_0\cap R_{\mathrm{reg}}$.
\end{theorem}

\begin{proof}
A good fan-out-two rank-one grammar is
\[
\begin{aligned}
 A_{\mathtt c}
   &\to(\mathtt a\mathtt c,\mathtt b)
     \mid(\mathtt a\xi_1^1\mathtt c,\xi_1^2\mathtt b)(A_{\mathtt c}),\\
 A_{\mathtt d}
   &\to(\mathtt a\mathtt d,\mathtt b\mathtt b)
     \mid(\mathtt a\xi_1^1\mathtt d,\xi_1^2\mathtt b\mathtt b)(A_{\mathtt d}),\\
 S&\to\mathtt c\mid\mathtt d
   \mid(\xi_1^1\mathtt c\xi_1^2)(A_{\mathtt c})
   \mid(\xi_1^1\mathtt d\xi_1^2)(A_{\mathtt d}).
\end{aligned}
\]
It generates exactly $L_\star$.  The intersection with the regular language
$\mathtt a^*\mathtt c^+\mathtt b^*$ is
$\{\mathtt a^n\mathtt c^{n+1}\mathtt b^n:n\ge0\}$, which is not
context-free.  For a short closure proof, introduce a fresh symbol
$\mathtt\#$ and the homomorphism fixing $\mathtt a,\mathtt b,\mathtt c$ and
sending $\mathtt\#$ to $\mathtt c$.  The inverse image of this branch,
intersected with
$\mathtt a^*\mathtt c^*\mathtt\#\mathtt b^*$, is
$\{\mathtt a^n\mathtt c^n\mathtt\#\mathtt b^n:n\ge0\}$; erasing
$\mathtt\#$ would make $\{\mathtt a^n\mathtt c^n\mathtt b^n:n\ge0\}$
context-free, a contradiction.  Thus the minimum fan-out is two, and
infinitude together with the displayed grammar gives minimum fan-out-two rule
rank one.

For typed substitutability put
\[
 Q_\star:=\mathtt a^*\mathtt c^+\mathtt b^*
          \cup\mathtt a^*\mathtt d^+\mathtt b^*,
 \qquad
 \Psi(w):=(|w|_{\mathtt a},|w|_{\mathtt b},
            |w|_{\mathtt c},|w|_{\mathtt d}).
\]
The strictly $2$-local language $Q_\star$ is boundary typed by
Lemma~\ref{lem:sl2-boundary-control}.  Let equal-$h_\star$-type tuples
$\tuple x,\tuple y$ share an accepting context $E_0$, and set
\[
 \Delta:=\sum_i\Psi(y_i)-\sum_i\Psi(x_i).
\]
Their equal $g_\star$-types put $E_0[\tuple x],E_0[\tuple y]$ on the same
branch, so for some $q\in\mathbb Z$,
\[
 \Delta=q(1,1,1,0)\quad\text{or}\quad\Delta=q(1,2,0,1)
\]
according as that branch is $\mathtt c$ or $\mathtt d$.

If $F[\tuple x]\in L_\star$, boundary typing gives
$F[\tuple y]\in Q_\star$, and equal $g_\star$-types put the two fillings on
the same $F$-branch.  When $E_0$ and $F$ choose the same branch, the displayed
count difference preserves its affine relation.  If they choose different
branches, no tuple component can contain $\mathtt c$ or $\mathtt d$; hence the
two center coordinates of $\Delta$ vanish, forcing $q=0$.  In either case
$F[\tuple y]\in L_\star$, and symmetry proves
$L_\star\in\Cmcf{f}{h_\star}$.

Finally assume $L_\star=L_0\cap R_{\mathrm{reg}}$ with
$L_0\in\YSub{f}$ and $R_{\mathrm{reg}}$ regular.  Since
$\mathtt c,\mathtt d\in L_0$, arity-one substitutability gives
$\D_{L_0}(\mathtt c)=\D_{L_0}(\mathtt d)$.  Hence
$\mathtt a\mathtt d^2\mathtt b^2\in L_0$ yields
$\mathtt a\mathtt c^2\mathtt b^2\in L_0$, and together with
$\mathtt a\mathtt c^2\mathtt b\in L_0$ gives
\[
 \D_{L_0}(\mathtt c^2\mathtt b)
 =\D_{L_0}(\mathtt c^2\mathtt b^2).
\]
Thus $\mathtt a^m\mathtt c^{m+1}\mathtt b^{m+t}\in L_0$ for all
$m\ge1,t\ge0$.  For $m$ at least the state number of a DFA for
$R_{\mathrm{reg}}$, pumping the final $\mathtt b^m$ block of the accepted
target word gives some $t>0$ with the same word in $R_{\mathrm{reg}}$, hence
in the intersection but not in $L_\star$, a contradiction.
\end{proof}

Thus the single target $L_\star$ rules out representation by a Yoshinaka
language followed by any finite-state filter, irrespective of how the filter
typing refines $h_\star$.

The witnesses $H_m$ are rank one and fall under
Theorem~\ref{thm:rankone-polydata}; by contrast,
Theorem~\ref{thm:intrinsic-ranktwo-witness} gives a boundary-typed
non-context-free target whose minimum rule rank remains two across all finite
fan-outs, placing the higher-rank boundary theorem genuinely beyond the
rank-one regime.


\section{Conclusion}
\label{sec:conclusion}

Finite-monoid typing gives a compositional bias for positive-data MCFG
learning: every fixed $(f,h)$ admits exact reconstruction and TxtEx
identification, with polynomial-time updates under a fixed rule-rank cap.
Fan-out is part of the bias: $L_{\mathrm{cd}}$ shows that increasing tested
arity can destroy substitutability beyond repair by any finite typing, while
Theorem~\ref{thm:every-fanout-strictness} gives typed rank-one witnesses of
minimum fan-out $m$ outside Yoshinaka's class.  The language $L_\star$
additionally excludes arbitrary target-dependent regular filtering.

For characteristic data, arbitrary fixed typing has a grammar-size polynomial
at rule rank one, while prefix--suffix-determined typing has a
grammar-size-and-rule-thickness polynomial at arbitrary fixed rank.
Theorem~\ref{thm:intrinsic-ranktwo-witness} supplies a boundary-typed
non-context-free target of minimum fan-out and rule rank two.  Full yield-type
refinement and empty-component elimination can nevertheless cause exponential
thickness growth; Appendix~\ref{app:normalization-details} records the latter
normalization details.

\begin{openproblem}[Original-thickness bound for arbitrary finite typing]
Fix $f,r_{\max}$ and a finite typing $h$.  Does every reduced good
rank-$\le r_{\max}$ presentation $G$ of a target in $\Cmcf{f}{h}$ admit
characteristic data of total positive length bounded by
\[
  \operatorname{poly}_{f,r_{\max},h}
  (\|G\|\max\{1,\theta_G\})\ ?
\]
The obstruction from Proposition~\ref{prop:typed-thickness-exponential} already
occurs at fan-out one and rule rank two, so this question is open already for
branching CFG presentations.  One possible route to a positive answer would be
to select a polynomially controlled sufficient typed subpresentation; another
would be to replace the target presentation by an equivalent one whose rule
children have unique componentwise $h$-types.  Either route would be
sufficient, but neither is known to be necessary.
\end{openproblem}



Other directions include adaptive learning over several fan-out biases,
learning useful typings without an a priori codomain bound, characterizing the
rank-one subfamily inside the typed classes, understanding the relation with
full-multicontext congruential MCFGs, and extending the reconstruction to
PMCFG-style copying.

\section*{Acknowledgments}
The author thanks Makoto Kanazawa and Ryo Yoshinaka for valuable discussions
and encouragement.

\section*{Declaration of competing interest}
The author declares that he has no known competing financial interests or
personal relationships that could have appeared to influence the work reported
in this paper.

\section*{Funding}
This research received no specific grant from any funding agency in the public,
commercial, or not-for-profit sectors.

\section*{Data availability}
No data were used for the research described in this article.

\section*{Declaration of generative AI and AI-assisted technologies in the writing process}
During the preparation of this manuscript, the author used ChatGPT and Claude
to assist with language editing, consistency checks, and critical review of the
draft.  The author reviewed and edited all resulting material and takes full
responsibility for the content of the article.

\appendix
\section{Details of the Good Normalization}
\label{app:normalization-details}

The following bookkeeping records the fixed-parameter bounds used in
Proposition~\ref{prop:good-normalization}; semantic correctness of the
nondeleting and nonmerging conversions is standard
\cite{SekiEtAl1991,Kanazawa2019Ogden,Yoshinaka2011}.

\paragraph{Nondeleting conversion and component permutations.}
After trimming, use the nondeleting step from the proof of Seki et al.'s
Lemma~2.2 in Kanazawa's rank-preserving form
\cite[proof of Lemma~2.2]{SekiEtAl1991}\cite[Section~2.1]{Kanazawa2019Ogden}.
Retained-component subsets give at most $O_f(1)$ decorated copies per
nonterminal and $O_{f,r_{\max}}(1)$ projected copies per rule; deleting an
empty child cannot increase rank.  For nonpermutation, decorate a fan-out-$d$
nonterminal by $\sigma\in\mathfrak S_d$.  The parent permutation determines
the child permutations, so there are at most $f!$ copies, no child is
duplicated, and rank is preserved.  Denote the result by $G_{\mathrm{nd}}$.

\paragraph{Removing empty tuple components.}
For a fan-out-$d$ nonterminal $A$, let
\[
 \mathsf{Pat}(A):=
 \{I\subseteq\{1,\ldots,d\}:\text{$I$ is the nonempty-component pattern of
 some }\tuple w\in L_A(G_{\mathrm{nd}})\}.
\]
Finite saturation computes these sets.  For realizable child patterns, replace
variables from absent components by $\eps$, delete empty parent components,
and omit all-empty children.  There are at most $2^f$ patterns per nonterminal
and $2^{fr_{\max}}$ choices per rule.  Introduce $A^I$ for each nonempty
pattern $I$.  Retained variables still occur once and in order, so
nondeletion/nonpermutation, fan-out, and rank are preserved and the result is
$\lambda$-free.

Induction on derivations gives
$L_{A^I}$ exactly as the projection of source tuples having pattern $I$;
all-empty omitted children are restored by their certified witnesses.  Since
the start symbol has fan-out one, $S^{\{1\}}$ generates precisely the nonempty
part of the original language.  Projection shortens templates, giving
$O_{f,r_{\max}}(\|G_{\mathrm{nd}}\|)$ size.

\paragraph{Nonmerging.}
Apply Yoshinaka's nonmerging conversion
\cite[Section~2.3, Lemma~1]{Yoshinaka2011}.  Each decorated nonterminal has
the full-image property used below, and adjacent-component concatenation
preserves $\lambda$-freeness, nonpermutation, fan-out, and rule rank.
For fixed fan-out and rank Yoshinaka's combined good-form conversion has
$O(\|G\|)$ size~\cite[Section~2.3, after Example~2]{Yoshinaka2011}; trimming
can only decrease all bounds.  This proves
Proposition~\ref{prop:good-normalization}.

\subsection*{Thickness transfer for $\lambda$-free presentations}

\begin{proposition}[Thickness under good-form conversion]
\label{prop:structural-good-thickness}
Fix $f$ and $r_{\max}$.  Let $G_{\mathrm{in}}$ be a reduced $\lambda$-free
$f$-MCFG of maximum rule rank at most $r_{\max}$, possibly with deleting
rules, and suppose $L(G_{\mathrm{in}})\cap\Sigma^+\neq\emptyset$.  The good
normalization of Proposition~\ref{prop:good-normalization} yields an equivalent
reduced good grammar $G$ with the same fan-out and rule-rank bounds such that
\[
  \|G\|=O_{f,r_{\max}}(\|G_{\mathrm{in}}\|),\qquad
  \tau_G\le \tau_{G_{\mathrm{in}}},\qquad
  \theta_G=O_{f,r_{\max}}\!\bigl(
     \|G_{\mathrm{in}}\|\max\{1,\tau_{G_{\mathrm{in}}}\}\bigr).
\]
\end{proposition}

\begin{proof}
For $\lambda$-free $G_{\mathrm{in}}$, every derived component is nonempty.
The nondeleting conversion therefore needs no later empty-pattern split, and
each retained decoration $A^I$ generates exactly the projection of
$L_A(G_{\mathrm{in}})$ onto $I$.  Hence its minimum total yield cannot exceed
that of $A$.  Component permutations preserve total length, and Yoshinaka's
nonmerging copies are exact adjacent-component images
\cite[Section~2.3, Lemma~1]{Yoshinaka2011}, so
$\tau_G\le\tau_{G_{\mathrm{in}}}$.  All decorations have fixed-parameter
linear overhead, giving
$\|G\|=O_{f,r_{\max}}(\|G_{\mathrm{in}}\|)$; then
Lemma~\ref{lem:thickness-equivalence} yields the displayed bound on
$\theta_G$.
\end{proof}

Thus deleting rules do not by themselves obstruct polynomial thickness
transfer from a reduced $\lambda$-free presentation; genuine empty components
are the problematic step.

\begin{proposition}[Empty-component elimination can blow up thickness]
\label{prop:empty-component-blowup}
For every $n\ge1$ there is a reduced fan-out-two, rule-rank-two MCFG $E_n$ of
size $O(n)$ with $\tau_{E_n}\le2$ and $\theta_{E_n}=O(n)$ such that the
empty-component elimination stage of Proposition~\ref{prop:good-normalization}
contains a retained nonterminal and a retained rule of thickness at least
$2^n$.
\end{proposition}

\begin{proof}
Let $E_n$ have fan-out-one $S$, fan-out-two $A_0,\ldots,A_n$, and
\[
\begin{aligned}
 A_0&\to(\mathtt a,\eps),\\
 A_i&\to(\eps,\mathtt c)\mid
      (\xi_1^1\xi_2^1,\xi_1^2\xi_2^2)(A_{i-1},A_{i-1})
      \quad(1\le i\le n),\\
 S&\to(\xi_1^1\mathtt d\,\xi_1^2)(A_n).
\end{aligned}
\]
It is reduced, nondeleting, nonpermuting, and nonmerging.  Short side
derivations give $\tau_{E_n}\le2$, and placing any chosen rule on one branch
of the binary spine while completing siblings minimally gives
$\theta_{E_n}=O(n)$.

If an $A_i$-tuple has nonempty first component and empty second component, the
rule $(\eps,\mathtt c)$ cannot occur.  Inductively the shortest such tuple is
$(\mathtt a^{2^i},\eps)$.  Thus pattern $\{1\}$ for $A_n$ survives
empty-component elimination, and the projected start rule
\[
 S\to\xi_1^1\mathtt d\,(A_n^{\{1\}})
\]
can succeed only with a child of length at least $2^n$.  Hence the retained
nonterminal and rule have thickness at least $2^n$.
\end{proof}

This is the same obstruction exhibited by
Proposition~\ref{prop:typed-thickness-exponential}: finite state refinement can
force exponentially long shortest witnesses inside one refined state.

\small
\begin{thebibliography}{99}

\bibitem{Gold1967}
E.~M.~Gold.
\newblock Language identification in the limit.
\newblock {\em Information and Control}, 10(5):447--474, 1967.

\bibitem{Angluin1980}
D.~Angluin.
\newblock Inductive inference of formal languages from positive data.
\newblock {\em Information and Control}, 45(2):117--135, 1980.

\bibitem{ClarkEyraud2007}
A.~Clark and R.~Eyraud.
\newblock Polynomial identification in the limit of substitutable context-free
languages.
\newblock {\em Journal of Machine Learning Research}, 8:1725--1745, 2007.

\bibitem{deLaHiguera1997}
C.~de~la~Higuera.
\newblock Characteristic sets for polynomial grammatical inference.
\newblock {\em Machine Learning}, 27:125--138, 1997.

\bibitem{CosteTyping2004}
F.~Coste, D.~Fredouille, C.~Kermorvant, and C.~de~la~Higuera.
\newblock Introducing domain and typing bias in automata inference.
\newblock In \emph{Grammatical Inference: Algorithms and Applications -- ICGI 2004},
LNCS 3264, pp.~115--126, Springer, 2004.

\bibitem{Takada1988}
Y.~Takada.
\newblock Grammatical inference for even linear languages based on control sets.
\newblock \emph{Information Processing Letters}, 28(4):193--199, 1988.
\newblock doi:10.1016/0020-0190(88)90208-6.

\bibitem{Takada1995}
Y.~Takada.
\newblock Learning formal languages based on control sets.
\newblock In K.~P. Jantke and S.~Lange (eds.), \emph{Algorithmic Learning
for Knowledge-Based Systems}, LNCS 961, pp.~316--339, Springer, 1995.

\bibitem{Clark2013}
A.~Clark.
\newblock Learning trees from strings: A strong learning algorithm for some context-free grammars.
\newblock \emph{Journal of Machine Learning Research}, 14:3537--3559, 2013.

\bibitem{CosteGaretNicolas2012}
F.~Coste, G.~Garet, and J.~Nicolas.
\newblock Local substitutability for sequence generalization.
\newblock In {\em Proceedings of the 11th International Conference on
Grammatical Inference}, JMLR Workshop and Conference Proceedings 21,
pp.~97--111, 2012.

\bibitem{EyraudHeinzYoshinaka2016}
R.~Eyraud, J.~Heinz, and R.~Yoshinaka.
\newblock Efficiency in the identification in the limit learning paradigm.
\newblock In J.~Heinz and J.~M.~Sempere (eds.), {\em Topics in Grammatical
Inference}, pp.~25--46. Springer, 2016.
\newblock doi:10.1007/978-3-662-48395-4\_2.

\bibitem{CosteNicolas2019}
F.~Coste and J.~Nicolas.
\newblock Learning local substitutable context-free languages from positive
examples in polynomial time and data by reduction.
\newblock In {\em Proceedings of the 14th International Conference on
Grammatical Inference}, Proceedings of Machine Learning Research 93,
pp.~155--168, 2019.

\bibitem{Yoshinaka2008}
R.~Yoshinaka.
\newblock Identification in the limit of $k,l$-substitutable context-free
languages.
\newblock In {\em Grammatical Inference: Algorithms and Applications},
LNCS 5278, pp.~266--279. Springer, 2008.

\bibitem{Yoshinaka2009}
R.~Yoshinaka.
\newblock Learning mildly context-sensitive languages with multidimensional
 substitutability from positive data.
\newblock In {\em Algorithmic Learning Theory}, LNCS 5809, pp.~278--292.
 Springer, 2009.
\newblock doi:10.1007/978-3-642-04414-4\_24.

\bibitem{Yoshinaka2011}
R.~Yoshinaka.
\newblock Efficient learning of multiple context-free languages with
multidimensional substitutability from positive data.
\newblock {\em Theoretical Computer Science}, 412(19):1821--1831, 2011.
\newblock doi:10.1016/j.tcs.2010.12.058.

\bibitem{Yoshinaka2010}
R.~Yoshinaka.
\newblock Polynomial-time identification of multiple context-free languages
from positive data and membership queries.
\newblock In {\em Grammatical Inference: Theoretical Results and Applications},
LNCS 6339, pp.~230--244. Springer, 2010.
\newblock doi:10.1007/978-3-642-15488-1\_19.

\bibitem{ClarkYoshinaka2016}
A.~Clark and R.~Yoshinaka.
\newblock Distributional learning of context-free and multiple context-free grammars.
\newblock In J.~Heinz and J.~M.~Sempere (eds.), \emph{Topics in Grammatical
Inference}, pp.~143--172. Springer, 2016.
\newblock doi:10.1007/978-3-662-48395-4\_6.

\bibitem{ClarkYoshinaka2014}
A.~Clark and R.~Yoshinaka.
\newblock Distributional learning of parallel multiple context-free grammars.
\newblock {\em Machine Learning}, 96(1--2):5--31, 2014.
\newblock doi:10.1007/s10994-013-5403-2.

\bibitem{OatesEtAl2006}
T.~Oates, T.~Armstrong, L.~Becerra-Bonache, and M.~Atamas.
\newblock Inferring grammars for mildly context-sensitive languages in polynomial-time.
\newblock In {\em Grammatical Inference: Algorithms and Applications -- ICGI 2006},
LNCS 4201, pp.~137--147. Springer, 2006.

\bibitem{KanazawaYoshinaka2023}
M.~Kanazawa and R.~Yoshinaka.
\newblock Extending distributional learning from positive data and membership queries.
\newblock In {\em Proceedings of the 16th International Conference on Grammatical Inference},
Proceedings of Machine Learning Research 217, pp.~8--22, 2023.
\bibitem{Gallot2021}
P.~D.~Gallot.
\newblock \emph{Safety of transformations of data trees: Tree transducer
theory applied to a verification problem on shell scripts}.
\newblock PhD thesis, Universit\'e de Lille, 2021.
\newblock
\url{https://pepite-depot.univ-lille.fr/LIBRE/EDMADIS/2021/2021LILUB030.pdf}.

\bibitem{BishopEtAl2026EDT0L}
A.~Bishop, M.~Elder, A.~Evetts, P.~Gallot, and A.~Levine.
\newblock On groups with EDT0L word problem.
\newblock \emph{International Journal of Algebra and Computation},
36(4):425--486, 2026.
\newblock doi:10.1142/S0218196726500165.

\bibitem{RambowSatta1999}
O.~Rambow and G.~Satta.
\newblock Independent parallelism in finite copying parallel rewriting systems.
\newblock {\em Theoretical Computer Science}, 223(1--2):87--120, 1999.
\newblock doi:10.1016/S0304-3975(97)00190-4.

\bibitem{kuriyamaCFG}
T.~Kuriyama.
\newblock Distributional learning of context-free languages under fixed finite-monoid typing.
\newblock Revised manuscript under review at \emph{Theoretical Computer Science}, 2026.
\newblock Public predecessor: arXiv:1409.6247v4 [cs.FL].

\bibitem{McNaughtonPapert1971}
R.~McNaughton and S.~Papert.
\newblock \emph{Counter-Free Automata}.
\newblock MIT Press, Cambridge, MA, 1971.

\bibitem{Pin1986}
J.-E. Pin.
\newblock {\em Varieties of Formal Languages}.
\newblock North Oxford Academic, London, and Plenum, New York, 1986.

\bibitem{SekiEtAl1991}
H.~Seki, T.~Matsumura, M.~Fujii, and T.~Kasami.
\newblock On multiple context-free grammars.
\newblock {\em Theoretical Computer Science}, 88(2):191--229, 1991.
\newblock doi:10.1016/0304-3975(91)90374-B.

\bibitem{Kallmeyer2010}
L.~Kallmeyer.
\newblock {\em Parsing Beyond Context-Free Grammars}.
\newblock Springer, Berlin, Heidelberg, 2010.
\newblock doi:10.1007/978-3-642-14846-0.

\bibitem{Kanazawa2019Ogden}
M.~Kanazawa.
\newblock Ogden's lemma, multiple context-free grammars, and the control language hierarchy.
\newblock {\em Information and Computation}, 269:104449, 2019.
\newblock doi:10.1016/j.ic.2019.104449.

\bibitem{GomezRodriguezEtAl2009}
C.~G\'omez-Rodr\'iguez, M.~Kuhlmann, G.~Satta, and D.~Weir.
\newblock Optimal reduction of rule length in linear context-free rewriting systems.
\newblock In {\em Proceedings of Human Language Technologies: The 2009 Annual
Conference of the North American Chapter of the Association for Computational
Linguistics}, pp.~539--547. Association for Computational Linguistics, 2009.


\end{thebibliography}
\normalsize

\end{document}